\chapimg{fig/op01.png}
\IfFileExists{extra/lead01.tex}{\chapterlead{\subhd{After this chapter you can} describe soil as a three-phase system and compute its porosity, void ratio and degrees of saturation; classify a soil by texture; express soil water as content, depth and suction and convert between them; derive the water balance of a root zone; write the van Genuchten retention and Mualem conductivity functions and evaluate them; and reduce the field data of a profile to the numbers (total available water, allowed depletion, interval) that an irrigation schedule needs.}
}{}
\chapter{Soil Water: Balance, Retention and Conductivity}\label{chap:01}
\markboth{Soil Water: Balance, Retention and Conductivity}{Soil Water: Balance, Retention and Conductivity}
\IfFileExists{extra/bg01.tex}{\section{Introduction: water, soil and agriculture}
Agriculture is the largest user of fresh water: about seventy per cent of the water withdrawn from rivers, lakes and aquifers goes to irrigation, and most of it does not come back to the stream it left. The water that a crop uses passes through one thin layer of the earth, the soil, and the behaviour of that layer decides how much of an applied millimetre is held for the roots, how much is lost to the air and how much drains out of reach. A designer of irrigation systems, a modeller of the soil--plant--atmosphere system and a farmer who reads a tensiometer all need the same small set of ideas: how much water the soil holds, how tightly, and how fast it moves.

This chapter gives those ideas in the order in which a textbook of soil physics gives them: first the soil as a mixture of solids, water and air; then its texture and the classes into which texture divides it; then the amount of water it holds and the ways of measuring that amount; then the energy with which the water is held, which decides where it moves; and last the characteristic constants (saturation, field capacity, wilting point) that turn all this into numbers for a schedule. The second half of the chapter, which carries over from the studies on which this book is based, derives the retention and conductivity functions that the whole of Daisy and of the SWAFL application use. Every definition is followed by a solved example; the chapter ends with a design project, exercises and questions.

\section{The soil as a three-phase system}
A soil is a porous body made of three phases: solid particles (mineral grains and organic matter), liquid water with dissolved salts, and gas (air with water vapour). If a sample of total volume $V_t$ is divided into the volume of solids $V_s$, of water $V_w$ and of air $V_a$, the volume of the pores is $V_v=V_w+V_a$ and $V_t=V_s+V_v$. The masses are $M_s$ (dry solids) and $M_w$ (water); the mass of air is neglected. The following quantities are defined from these six numbers.

\begin{equation}\rho_b=\frac{M_s}{V_t},\qquad \rho_s=\frac{M_s}{V_s},\qquad \phi=\frac{V_v}{V_t}=1-\frac{\rho_b}{\rho_s},\qquad e=\frac{V_v}{V_s}=\frac{\phi}{1-\phi}\label{eq:phase}\end{equation}


Here $\rho_b$ is the dry bulk density, $\rho_s$ the particle density (about 2.65 Mg\,m$^{-3}$ for mineral soils), $\phi$ the porosity and $e$ the void ratio. The water content is described in two ways. The gravimetric content $w=M_w/M_s$ is a mass ratio and is what a laboratory oven measures; the volumetric content $\theta=V_w/V_t$ is what a model uses, because it multiplies a depth. The two are related by $\theta=w\,\rho_b/\rho_w$, with $\rho_w=1$ Mg\,m$^{-3}$. The degree of saturation is $S=V_w/V_v=\theta/\phi$ and the air-filled porosity is $\phi_a=\phi-\theta$.

\exbox{Phase relations of a core}{An undisturbed core of volume $V_t=250$ cm$^3$ has a wet mass of 402 g and, after drying at 105\,$^\circ$C, a dry mass of 337.5 g. Take $\rho_s=2.65$ Mg\,m$^{-3}$. Find the dry bulk density, the porosity, the void ratio, the gravimetric and volumetric water contents, the degree of saturation and the air-filled porosity.}\sol Given: $V_t=250$ cm$^3$, $M_{wet}=402$ g, $M_s=337.5$ g. \subhd{Bulk density.} $\rho_b=337.5/250=1.35$ g\,cm$^{-3}$ \ans \subhd{Porosity and void ratio.} $\phi=1-1.35/2.65=0.4906$ and $e=0.4906/(1-0.4906)=0.963$ \ans \subhd{Water contents.} $M_w=402-337.5=64.5$ g, so $w=64.5/337.5=0.1911$. Then $\theta=w\rho_b=0.1911\times1.35=0.258$ cm$^3$\,cm$^{-3}$ \ans \subhd{Saturation and air.} $S=0.258/0.4906=0.526$ and $\phi_a=0.4906-0.258=0.233$ \ans So about half of the pore space holds water and a quarter of the soil volume is air.

\section{Soil texture and the classes of soils}
The mineral part of a soil is a mixture of particles of different sizes. By the USDA convention, \kw{sand} is the fraction between 2 and 0.05 mm, \kw{silt} between 0.05 and 0.002 mm and \kw{clay} below 0.002 mm; the proportions of the three (summing to 100 per cent) are the \kw{texture}. The texture triangle divides all possible mixtures into twelve classes. In irrigation the classes are often grouped in three types: \emph{coarse} soils (sand, loamy sand, sandy loam), which take in water fast, hold little and drain quickly; \emph{medium} soils (loam, silt loam, silt), which hold the most plant-available water; and \emph{fine} soils (clay loam, silty clay loam, silty clay, clay), which take in water slowly, hold it tightly and are prone to runoff and waterlogging.

The ten soils used throughout this book are the ones of the Daisy ensemble: S01 coarse loam, S02 coarse sand, S03 sand, S04 loamy sand, S05 sandy loam, S06 loam, S07 silt loam, S08 sandy clay loam, S09 clay loam and S10 silty clay loam. Table~\ref{tab:ten} lists their composition and van Genuchten parameters; the parameter $\alpha$ is, roughly, the inverse of the air-entry suction and $n$ describes the width of the pore-size distribution (a large $n$ is a uniform pore size, as in sand).

\begin{table}[!htbp]\centering\small\caption{The ten soils of the ensemble: composition (per cent by mass), van Genuchten parameters and plant-available water per centimetre of depth.}
\begin{tabular}{llrrrrrr}\toprule
\hd{Soil} & \hd{Class} & \hd{Sand \%} & \hd{Silt \%} & \hd{Clay \%} & \hd{$\alpha$ (cm$^{-1}$)} & \hd{$n$} & \hd{AWC (mm cm$^{-1}$)}\\\midrule
S01 & Coarse loam & 45 & 35 & 20 & 0.0961 & 1.353 & 1.5\\
S02 & Coarse sand & 88 & 8 & 4 & 0.0536 & 1.600 & 0.5\\
S03 & Sand & 92 & 5 & 3 & 0.1450 & 2.680 & 0.6\\
S04 & Loamy sand & 82 & 12 & 6 & 0.1240 & 2.280 & 0.9\\
S05 & Sandy loam & 65 & 25 & 10 & 0.0750 & 1.890 & 1.2\\
S06 & Loam & 40 & 40 & 20 & 0.0360 & 1.560 & 1.7\\
S07 & Silt loam & 20 & 65 & 15 & 0.0200 & 1.410 & 2.0\\
S08 & Sandy clay loam & 55 & 15 & 30 & 0.0590 & 1.480 & 1.4\\
S09 & Clay loam & 30 & 35 & 35 & 0.0190 & 1.310 & 1.6\\
S10 & Silty clay loam & 10 & 55 & 35 & 0.0100 & 1.230 & 1.7\\\bottomrule\end{tabular}\end{table}

\section{The water content of a soil and how it is measured}
The amount of water in a soil can be measured directly or inferred. The methods fall into four types. \kw{Gravimetric sampling} weighs a sample before and after oven drying; it is the reference method but destroys the sample and gives one point in time. \kw{Neutron scattering} counts the slow neutrons that return to a probe after hydrogen nuclei in the soil have thermalised them; it integrates a sphere of 15--25 cm radius. \kw{Dielectric methods} (time-domain reflectometry, capacitance and frequency-domain sensors) measure the dielectric permittivity of the soil, which is about 80 for water and 3--5 for dry soil, and convert it to $\theta$ by the Topp relation $\theta=-5.3\times10^{-2}+2.92\times10^{-2}\varepsilon-5.5\times10^{-4}\varepsilon^2+4.3\times10^{-6}\varepsilon^3$. \kw{Remote sensing} uses the microwave and thermal signal of the surface, and reaches only the top few centimetres.

A water content is turned into a \emph{depth} by multiplying by the thickness: a layer of thickness $z$ holds the depth $d=\theta z$ of water, in the same unit as $z$. Depths of water are the language of irrigation, because rain, evaporation and irrigation are all written in millimetres; one millimetre over one hectare is 10 m$^3$.

\exbox{Water in a profile}{A profile is sampled in four layers. Gravimetric contents on the sampling day: 0--20 cm, $w=0.140$ ($\rho_b=1.30$); 20--40 cm, $w=0.165$ ($\rho_b=1.38$); 40--80 cm, $w=0.190$ ($\rho_b=1.45$); 80--100 cm, $w=0.205$ ($\rho_b=1.50$ g\,cm$^{-3}$). Find the volumetric content of each layer, the depth of water held in each, and in the whole profile.}\sol For each layer $\theta=w\rho_b$ and $d=\theta z$. layer 0--20: $\theta=0.140\times1.30=0.182$, $d=0.182\times200=36.4$ mm; layer 20--40: $\theta=0.165\times1.38=0.228$, $d=0.228\times200=45.5$ mm; layer 40--80: $\theta=0.190\times1.45=0.275$, $d=0.275\times400=110.2$ mm; layer 80--100: $\theta=0.205\times1.50=0.307$, $d=0.307\times200=61.5$ mm. The sum is $d=253.6$ mm \ans

\section{Soil water potential}
Water in a soil has energy, and it moves from where its energy is high to where it is low. The \kw{total potential} $\Psi_t$ is the work needed to move a unit quantity of pure free water from a reference pool to the point in the soil; per unit \emph{weight} it is a length, the \kw{hydraulic head} $H$. Four components are distinguished: the \emph{matric} (capillary and adsorptive) potential $h$, always negative in an unsaturated soil; the \emph{gravitational} potential $z$, the height above a datum; the \emph{osmotic} potential, due to dissolved salts, which acts on root uptake but not on flow in the soil; and the \emph{pressure} potential, positive below a water table. Without salt, $H=h+z$ and flow is driven by the gradient of $H$ (Chapter~2). It is convenient to speak of the \kw{suction} $|h|$ (in cm of water, in kPa with 1 kPa $\approx 10.2$ cm, or as $\mathrm{pF}=\log_{10}|h|$).

\exbox{Heads and flow direction}{A tensiometer at 30 cm depth reads a suction of 250 cm of water, and one at 60 cm reads 210 cm (a datum at the surface, $z$ positive upward). Find the total head at each depth and the direction and size of the flow if the mean conductivity between them is $K=0.8$ cm\,d$^{-1}$.}\sol Heads: $H_{30}=h+z=-250+(-30)=-280$ cm and $H_{60}=-210+(-60)=-270$ cm. \subhd{Direction.} Water moves from high head to low head, that is from 60 cm ($-270$) to 30 cm ($-280$): it moves \emph{upward}. \subhd{Flux.} Darcy's law gives $q=-K\,\Delta H/\Delta z=-0.8\times[(-280)-(-270)]/(30-60)$ with $z$ taken downward positive: the gradient is $10/30=0.333$, so $q=0.8\times0.333=0.267$ cm\,d$^{-1}$ upward \ans This is 2.7 mm per day, comparable to the evaporative demand: the water table is not needed to explain it.

\section{Characteristic soil water constants}
Because $\theta$ and $h$ are linked by the retention curve (later in this chapter), a few suctions have been given names. \kw{Saturation} ($h=0$) is the content at which all pores are filled. \kw{Field capacity} is the content that remains when drainage has become negligible, two or three days after a soaking, and is placed at $h=100$ cm (pF 2.0) in this book (some authors use 330 cm for fine soils). The \kw{permanent wilting point} is the content at which a crop cannot recover its turgor, placed at 15\,849 cm (pF 4.2, 1.55 MPa). The difference is the \kw{plant-available water}, and over the depth $Z_r$ of the roots the \kw{total available water} is

\begin{equation}\mathrm{TAW}=1000\,(\theta_{FC}-\theta_{WP})\,Z_r\qquad(\text{mm, with }Z_r\text{ in m}).\label{eq:taw}\end{equation}
\par\noindent where ${\rm TAW}$ is the total available water (mm), $\theta_{FC}$ and $\theta_{WP}$ the water contents at field capacity and wilting point (m$^3$\,m$^{-3}$), $Z_r$ the rooting depth (m) and 1000 converts m to mm.



A crop is allowed to use a fraction $p$ of the TAW without stress; the \kw{readily available water} is $\mathrm{RAW}=p\,\mathrm{TAW}$, and the depth that has been used is the \kw{depletion} $D_r$. Irrigation is called when $D_r$ reaches RAW and refills it to field capacity, so the net irrigation depth is $Z_n=D_r$ and the interval is $Z_n/\mathrm{ET}_c$.

\exbox{Soil water constants of the loam}{For the loam S06 ($\alpha=0.036$ cm$^{-1}$, $n=1.56$, $\theta_s=0.468$, $\theta_r=0.02$) find the water contents at field capacity (100 cm) and wilting point (15\,849 cm), the plant-available water per metre of soil, and the TAW and RAW ($p=0.5$) for a root depth of 0.8 m.}\sol With $m=1-1/n=0.359$. \subhd{Field capacity.} $(\alpha h)^n=(3.6)^{1.56}=7.38$, so $S_e=(1+7.38)^{-0.359}=0.466$ and $\theta_{FC}=0.02+0.448\times0.466=0.229$. \subhd{Wilting point.} $(\alpha h)^n=(570.6)^{1.56}=19945$ so $S_e=0.0286$ and $\theta_{WP}=0.033$. \subhd{Available water.} $\theta_{FC}-\theta_{WP}=0.196$, that is 196 mm per metre \ans \subhd{TAW and RAW.} $\mathrm{TAW}=1000\times0.196\times0.8=157$ mm and $\mathrm{RAW}=0.5\times157=78$ mm \ans

}{}
\banner{Water, soil and the problem of saving}
\section{Why a saved millimetre is not a gained millimetre}
Irrigation takes most of the freshwater that people withdraw, so a better use of each millimetre is a first-order goal. Evaporation is the largest single outflow from land that is not saturated, so any measure that lowers evaporation, such as a mulch, seems to promise water that is saved. The promise rests on an assumption: that water not evaporated stays available. But water is conserved. If it does not evaporate it either remains in the root zone, where the crop can use it, or it drains below the roots, where the crop cannot. Which of the two happens depends on the soil, the climate and the management. This book gives the equations that decide it.

\fig{p1_framework}{0.85}{The framework of the work: a designed ensemble of Daisy runs, the closed water balance and the four lines of analysis (reach, fate, memory, audit).}

\section{The water balance of a root zone}
Over a period, a root zone with rainfall $P$ and irrigation $I$ satisfies
\begin{equation}\label{eq:bal}\boxed{P + I = T + E_s + E_o + D + \Delta S + Q ,}\end{equation}
where $T$ is transpiration, $E_s$ soil evaporation, $E_o$ the other evaporation (intercepted water, ponds, litter and snow), $D$ drainage below the root zone, $\Delta S$ the change of stored water and $Q$ runoff. All terms are depths of water, in mm, per year if the period is a year. The sum $E_s+E_o+T$ is the evapotranspiration $ET$. Equation~(\ref{eq:bal}) is an accounting identity, and in a model it can be checked, because every term is written to a file.

\exbox{Closing a balance}{In one run of the ensemble the rainfall is $P=620$ mm\,yr$^{-1}$ and the irrigation $I=164$ mm\,yr$^{-1}$. The evapotranspiration is $ET=560$, drainage $D=80$ and runoff $Q=4$. The change of storage is then
$\Delta S = P+I-ET-D-Q = 620+164-560-80-4 = 140$ mm over the period. If a user had instead reported a storage change of $20$ mm, the balance would be open by $120$ mm, and one would ask which flux has been left out. Chapter~\ref{ch:audit} shows that in Daisy output such an open term exists and has a name.}



\section{Notation}
\begin{center}\small\begin{tabular}{@{}lll@{}}\toprule
\hd{Symbol} & \hd{Meaning} & \hd{Unit}\\\midrule
$h$ & suction head & cm\\
$\theta,\ \Se$ & water content; effective saturation & $-$\\
$\alpha,\ n,\ m$ & van Genuchten parameters, $m=1-1/n$ & cm$^{-1}$, $-$, $-$\\
$\Ks,\ K,\ \Kr$ & saturated, unsaturated and relative conductivity & cm\,d$^{-1}$, $-$\\
$\ell$ & pore-connectivity exponent (0.5) & $-$\\
$\hc,\ h_t$ & reach of a trigger; trigger value & cm\\
$\kappa$ & relative conductivity at the reach & $-$\\
$\eta_X$ & marginal fate of irrigation in flux $X$ & $-$\\
$d,\ d_D$ & dose per event; dose scale of the drained share & mm\\
$\phi_x$ & memory fate of flux $x$ & $-$\\
$\bar R$ & released storage & mm\,yr$^{-1}$\\
$R$ & unattributed evaporative term & mm\,yr$^{-1}$\\\bottomrule
\end{tabular}\end{center}

\section{Three ways of knowing, and which this book uses}
A \emph{measurement} tells you what happened in a field. A \emph{model} tells you what follows from stated equations. A \emph{simulation ensemble} varies a model over many soils and conditions, so that the response of the model to each factor can be separated from the others. This book is about the third: the Daisy model is run 10,190 times, and the equations of the book are either derived from physics and checked against the ensemble, or fitted to the ensemble and tested on soils left out of the fit. The limits of this approach are stated in the last chapter.

\exbox{}{A root zone receives $P=480$ mm and $I=0$ in a rainfed year, with $T=210$, $E_s=170$, $E_o=60$, $Q=0$ and $D=55$. Find $\Delta S$.}\sol  $\Delta S=480-(210+170+60+0+55)=-15$ mm: the store lost 15 mm.

\exbox{}{Explain in two sentences why lowering $E_s$ by 50 mm does not necessarily raise $T$ by 50 mm.}\sol  Lowering $E_s$ leaves energy and water that may go to drainage, to storage or to transpiration; how much of each depends on the soil (Chapter 7).

\exbox{}{Using Equation~(\ref{eq:bal}), write the balance of the rainfed twin of an irrigated run and subtract the two. Show that dividing by $I$ gives five shares that sum to one.}\sol  Subtracting the two balances gives $I=(T-T_0)+(E_s-E_{s,0})+(E_o-E_{o,0})+(D-D_0)+(\Delta S-\Delta S_0+Q-Q_0)$; divide by $I$.
\banner{Soil water retention}
\section{Capillarity and the Young--Laplace law}
A soil is a porous medium. Water is held in the pores by surface tension, and the smaller the pore, the more strongly it holds. For a cylindrical capillary of radius $r$ and a contact angle $\gamma$, the pressure difference across the curved water surface is
\begin{equation}\label{eq:yl}\boxed{\Delta p = \frac{2\sigma\cos\gamma}{r},}\end{equation}
where $\sigma$ is the surface tension of water. Expressing the pressure as a suction head $h = \Delta p/(\rho_w g)$ gives
\begin{equation}\label{eq:hr}\boxed{h = \frac{2\sigma\cos\gamma}{\rho_w g\, r}\quad\Longrightarrow\quad h\propto \frac1r .}\end{equation}
\par\noindent where $h$ is the capillary rise or suction head (m, or cm), $\sigma$ the surface tension of water (0.0728 N\,m$^{-1}$ at 20$^\circ$C), $\gamma$ the contact angle (rad; about 0 for a wetting soil), $\rho_w$ the density of water (1000 kg\,m$^{-3}$), $g$ the gravitational acceleration (9.81 m\,s$^{-2}$) and $r$ the radius of the pore (m).

When the soil is drained to a suction $h$, pores with radius larger than $r(h)$ are empty and smaller pores are full. The retention curve is therefore a map of the pore-size distribution: the fraction of the pore space that is water-filled at suction $h$ equals the fraction of pores smaller than $r(h)$.

\exbox{The pore radius at 300 cm}{With $\sigma=0.0728$ N\,m$^{-1}$, $\gamma=0$, $\rho_w=998$ kg\,m$^{-3}$ and $g=9.81$ m\,s$^{-2}$, a suction of $h=300$ cm $=3$ m gives $r=2\sigma/(\rho_w g h)=2(0.0728)/(998\cdot 9.81\cdot 3)=4.96\times10^{-6}$ m, that is about 5 $\mu$m. At $h=10{,}000$ cm the radius is 0.15 $\mu$m. A trigger of 300 cm therefore says: ``irrigate when the pores larger than 5 $\mu$m are empty''.}



\begin{kfig}{The water balance of a root zone: rain $P$ and irrigation $I$ enter; transpiration $T$, evaporation $E$, drainage $D$ and runoff $Q$ leave; the store changes by $\Delta S$.}
\draw[soil] (0,0) rectangle (7,1.8);\draw[ln] (0,1.8)--(7,1.8);
\node[lab] at (3.5,0.9) {root zone: store $S$, change $\Delta S$};
\foreach \x/\l in {1.2/P,2.4/I}{\draw[vec] (\x,3.4)--(\x,1.9);\node[lab,above] at (\x,3.4) {$\l$};}
\foreach \x/\l in {4.4/T,5.6/E}{\draw[vec] (\x,1.9)--(\x,3.4);\node[lab,above] at (\x,3.4) {$\l$};}
\draw[vec] (3.5,0)--(3.5,-1.1);\node[lab,below] at (3.5,-1.1) {$D$};
\draw[vec] (7,1.4)--(8.1,1.4);\node[lab,right] at (8.1,1.4) {$Q$};
\draw[gd] (0,-0.5)--(7,-0.5);\node[lab,left] at (0,-0.5) {root depth};
\node[lab] at (3.5,-1.9) {$P+I=T+E+D+Q+\Delta S$};
\end{kfig}

\section{The effective saturation}
The effective saturation is the scaled water content
\begin{equation}\label{eq:Se}\boxed{\Se(h)=\frac{\theta(h)-\theta_r}{\theta_s-\theta_r},}\end{equation}
\par\noindent where $\Se$ is the effective saturation (dimensionless, 0 to 1), $\theta(h)$ the volumetric water content at suction $h$ (m$^3$\,m$^{-3}$), $\theta_r$ the residual water content and $\theta_s$ the saturated water content (both m$^3$\,m$^{-3}$).

which runs from 1 at saturation to 0 at the residual water content $\theta_r$.

\section{The van Genuchten curve}
Van Genuchten (1980) chose a smooth, closed form that follows the shape of measured curves:
\begin{equation}\label{eq:vg}\boxed{\Se(h)=\bigl[1+(\alpha h)^n\bigr]^{-m},\qquad m=1-\frac1n .}\end{equation}
\par\noindent where $h$ is the suction (cm), $\alpha$ the scale parameter (cm$^{-1}$, about the inverse of the air-entry suction), $n$ the shape parameter (dimensionless, $n>1$, larger for coarse and uniform soils) and $m$ the derived shape parameter (dimensionless).

The parameter $\alpha$ (cm$^{-1}$) sets the suction scale, roughly the inverse of the air-entry suction. The parameter $n$ sets the width of the pore-size distribution: a large $n$ means that most pores have nearly the same size, so the soil releases most of its water over a narrow band of suction; a small $n$ means a broad distribution and a gradual release. Inverting Equation~(\ref{eq:vg}),
\begin{equation}\label{eq:vginv}\boxed{h(\Se)=\frac1\alpha\Bigl(\Se^{-1/m}-1\Bigr)^{1/n}.}\end{equation}
\par\noindent where $h$ is the suction (cm) at the effective saturation $\Se$; $\alpha$, $n$ and $m$ are the van Genuchten parameters of Eq.~(\ref{eq:vg}).


\exbox{Effective saturation of sand and of silty clay loam}{Sand (S03) has $\alpha=0.145$ cm$^{-1}$ and $n=2.68$, so $m=1-1/2.68=0.627$. At $h=100$ cm, $\alpha h=14.5$ and $(\alpha h)^n=14.5^{2.68}=1.29\times10^{3}$. Then $\Se=(1+1290)^{-0.627}=0.011$: the sand is almost dry at 100 cm. Silty clay loam (S10) has $\alpha=0.010$ and $n=1.23$, $m=0.187$. At the same suction $\alpha h=1$, $(\alpha h)^n=1$ and $\Se=2^{-0.187}=0.878$: the soil is still 88 per cent saturated. Same suction, very different state of water.}



\fig{p4_curves}{0.85}{Retention and conductivity curves of the ten soils of the ensemble (from the study of the reach of a trigger).}

\begin{kfig}{Capillary rise in a tube of radius $r$: the meniscus (contact angle $\gamma$) holds the water column $h$ above the free surface, $h=2\sigma\cos\gamma/(\rho g r)$.}
\draw[wtr] (-2.2,0) rectangle (2.2,-1);\draw[ln] (-2.2,0)--(2.2,0);
\draw[ln] (-0.35,0)--(-0.35,3.2);\draw[ln] (0.35,0)--(0.35,3.2);
\fill[khcyan!25] (-0.35,0) rectangle (0.35,2.5);\draw[crv] (-0.35,2.5) to[bend right=40] (0.35,2.5);
\draw[<->,ln] (1.2,0)--(1.2,2.45);\node[lab,right] at (1.2,1.2) {$h$};
\draw[<->,ln] (-0.35,3.55)--(0.35,3.55);\node[lab,above] at (0,3.55) {$2r$};
\draw[vec] (0,2.0)--(0,1.1);\node[lab,left] at (-0.45,1.6) {$\rho g h$};
\draw[gd] (-1.4,2.5)--(1.1,2.5);\node[lab,right] at (2.3,0) {free water surface};
\node[lab] at (0.9,2.75) {$\gamma$};
\end{kfig}

\section{Where the parameters come from in Daisy}
\begin{daisybox}{soil parameters}
Daisy reads the van Genuchten parameters of each soil layer in its set-up file. In the ensemble, eight soils use the mean values of Carsel and Parrish (1988) and two use effective parameters. Table~\ref{tab:soils} lists them. The parameters are inputs. Daisy's retention function is then used inside the Richards equation, so every suction in the model logs follows Equation~(\ref{eq:vg}).
\end{daisybox}
\begin{table}[!htbp]\centering\small
\caption{The ten soils of the ensemble.}\label{tab:soils}
\begin{tabular}{llrrrr}\toprule
\hd{ID} & \hd{Texture} & \hd{$\alpha$ (cm$^{-1}$)} & \hd{$n$} & \hd{TAW (mm)} & \hd{Memory fate}\\\midrule
S01&Coarse loam&0.0961&1.353&130&0.029\\S02&Coarse sand&0.0536&1.600&105&0.265\\S03&Sand&0.1450&2.680&60&0.846\\S04&Loamy sand&0.1240&2.280&75&0.741\\S05&Sandy loam&0.0750&1.890&95&0.643\\S06&Loam&0.0360&1.560&150&0.289\\S07&Silt loam&0.0200&1.410&220&0.147\\S08&Sandy clay loam&0.0590&1.480&160&0.285\\S09&Clay loam&0.0190&1.310&175&0.087\\S10&Silty clay loam&0.0100&1.230&195&0.022\\\bottomrule\end{tabular}\end{table}

\exbox{}{Compute $\Se$ for the loam S06 ($\alpha=0.036$, $n=1.56$) at $h=50$, $300$ and $3000$ cm.}\sol  $\Se=0.638$ at 50 cm, $0.262$ at 300 cm and $0.073$ at 3000 cm.

\exbox{}{Show that Equation~(\ref{eq:vginv}) is the inverse of Equation~(\ref{eq:vg}).}\sol  Put $\Se=[1+(\alpha h)^n]^{-m}$; then $\Se^{-1/m}=1+(\alpha h)^n$, so $h=\alpha^{-1}(\Se^{-1/m}-1)^{1/n}$.

\exbox{}{Compute the radius of the largest water-filled pore at $h=1000$ cm and at $h=15{,}000$ cm (the permanent wilting point is often taken at 15,000 cm).}\sol  $r=0.1456\times10^{-3}/(998\cdot9.81\cdot h)$ with $h$ in metres: $1.49\ \mu$m at 1000 cm and $0.099\ \mu$m at 15,000 cm.

\exbox{}{Soil A has $n=2.5$ and soil B has $n=1.3$, with the same $\alpha$. Which one releases more water between 100 and 300 cm? Justify your answer with Equation~(\ref{eq:vg}).}\sol  Soil B (the broad one) releases more. With $\alpha=0.05$: $\Se(100)=0.089$ and $\Se(300)=0.017$ for $n=2.5$ (change 0.072); $0.601$ and $0.441$ for $n=1.3$ (change 0.160). The broad-pore soil has a flatter curve, so it gives up more water over this band.
\banner{Hydraulic conductivity}
\section{Why the conductivity falls}
Water flows through the pores that are water-filled. When the soil dries, the largest pores empty first, and these carry most of the flow, because the flow in a capillary goes as the fourth power of its radius (Poiseuille). Water must also find a connected, tortuous path. Both effects reduce the conductivity, and together they make it fall by many orders of magnitude.

\begin{kfig}{Retention curve of a loam (van Genuchten, $\alpha=0.036$ cm$^{-1}$, $n=1.56$): effective saturation against the logarithm of suction; field capacity at 100 cm and wilting point at 15\,849 cm.}
\draw[axis] (0,0)--(7.4,0);\draw[axis] (0,0)--(0,4.6);
\node[lab,below] at (3.7,-0.55) {Suction $h$ (cm, log scale)};\node[lab,rotate=90] at (-0.9,2.3) {Effective saturation $S_e$};
\foreach \x/\t in {0/1,1.5/10,3/100,4.5/1000,6/10000}{\draw[ln] (\x,0)--(\x,-0.1);\node[lab,below] at (\x,-0.1) {\t};}
\foreach \y/\t in {0/0,1.9/0.5,3.8/1}{\draw[ln] (0,\y)--(-0.1,\y);\node[lab,left] at (-0.1,\y) {\t};}
\draw[crv,smooth] plot coordinates {(-0.450,3.797) (-0.337,3.797) (-0.225,3.796) (-0.112,3.794) (0.000,3.792) (0.113,3.790) (0.225,3.787) (0.338,3.783) (0.450,3.778) (0.563,3.771) (0.675,3.762) (0.787,3.751) (0.900,3.736) (1.013,3.717) (1.125,3.693) (1.237,3.662) (1.350,3.623) (1.462,3.575) (1.575,3.514) (1.688,3.441) (1.800,3.353) (1.912,3.250) (2.025,3.132) (2.138,2.998) (2.250,2.852) (2.362,2.696) (2.475,2.532) (2.587,2.365) (2.700,2.198) (2.812,2.033) (2.925,1.874) (3.038,1.722) (3.150,1.578) (3.263,1.443) (3.375,1.318) (3.488,1.202) (3.600,1.095) (3.713,0.996) (3.825,0.906) (3.938,0.824) (4.050,0.749) (4.163,0.681) (4.275,0.618) (4.388,0.562) (4.500,0.510) (4.613,0.463) (4.725,0.421) (4.838,0.382) (4.950,0.347) (5.062,0.315) (5.175,0.286) (5.288,0.260) (5.400,0.236) (5.513,0.214) (5.625,0.194) (5.738,0.176) (5.850,0.160) (5.963,0.145) (6.075,0.132) (6.188,0.120) (6.300,0.109)};
\draw[gd] (3,0)--(3,1.52);\node[lab,above] at (3,0.05) {FC};
\draw[gd] (6.3,0)--(6.3,1);\node[lab,above] at (6.3,1.05) {WP};
\end{kfig}

\section{Mualem's model}
Mualem (1976) wrote the relative conductivity $\Kr=K/\Ks$ of a bundle of capillaries as the product of a tortuosity--connectivity term $\Se^{\ell}$ and the square of a normalised integral of the pore radius over the filled pores. Since $r\propto 1/h$ (Equation~\ref{eq:hr}),
\begin{equation}\label{eq:mualem}\boxed{\Kr(\Se)=\Se^{\ell}\left[\frac{\int_0^{\Se} h(x)^{-1}\,dx}{\int_0^{1} h(x)^{-1}\,dx}\right]^{2}.}\end{equation}
\par\noindent where $\Kr=K/K_s$ is the relative hydraulic conductivity (dimensionless), $K$ the conductivity at the given saturation and $K_s$ the saturated conductivity (both cm\,d$^{-1}$), $\ell$ the pore-connectivity parameter (dimensionless, 0.5 by default), $x$ a dummy variable of integration (an effective saturation) and $h(x)$ the suction at $x$ (cm).


\begin{kfig}{Relative conductivity of two soils against suction: the conductivity falls by orders of magnitude as the soil dries, and falls sooner in a coarse soil.}
\draw[axis] (0,0)--(7.4,0);\draw[axis] (0,0)--(0,4.4);
\node[lab,below] at (3.7,-0.55) {Suction $h$ (cm, log scale)};\node[lab,rotate=90] at (-0.9,2.2) {$\log_{10}K_r$};
\foreach \x/\t in {0/1,1.5/10,3/100,4.5/1000,6/10000}{\draw[ln] (\x,0)--(\x,-0.1);\node[lab,below] at (\x,-0.1) {\t};}
\draw[crv,smooth] plot coordinates {(-0.300,3.992) (-0.190,3.989) (-0.080,3.986) (0.030,3.981) (0.140,3.974) (0.250,3.966) (0.360,3.954) (0.470,3.938) (0.580,3.916) (0.690,3.887) (0.800,3.847) (0.910,3.794) (1.020,3.724) (1.130,3.634) (1.240,3.521) (1.350,3.384) (1.460,3.226) (1.570,3.049) (1.680,2.856) (1.790,2.653) (1.900,2.441) (2.010,2.224) (2.120,2.004) (2.230,1.781) (2.340,1.556) (2.450,1.330) (2.560,1.104) (2.670,0.878) (2.780,0.651) (2.890,0.424) (3.000,0.250) (3.110,0.250) (3.220,0.250) (3.330,0.250) (3.440,0.250) (3.550,0.250) (3.660,0.250) (3.770,0.250) (3.880,0.250) (3.990,0.250) (4.100,0.250) (4.210,0.250) (4.320,0.250) (4.430,0.250) (4.540,0.250) (4.650,0.250) (4.760,0.250) (4.870,0.250) (4.980,0.250) (5.090,0.250) (5.200,0.250) (5.310,0.250) (5.420,0.250) (5.530,0.250) (5.640,0.250) (5.750,0.250) (5.860,0.250) (5.970,0.250) (6.080,0.250) (6.190,0.250) (6.300,0.250)};
\draw[ln,smooth] plot coordinates {(-0.300,3.921) (-0.190,3.915) (-0.080,3.908) (0.030,3.900) (0.140,3.892) (0.250,3.883) (0.360,3.874) (0.470,3.863) (0.580,3.852) (0.690,3.839) (0.800,3.825) (0.910,3.810) (1.020,3.793) (1.130,3.774) (1.240,3.753) (1.350,3.730) (1.460,3.704) (1.570,3.675) (1.680,3.643) (1.790,3.608) (1.900,3.569) (2.010,3.525) (2.120,3.476) (2.230,3.423) (2.340,3.365) (2.450,3.301) (2.560,3.232) (2.670,3.158) (2.780,3.079) (2.890,2.995) (3.000,2.907) (3.110,2.815) (3.220,2.720) (3.330,2.622) (3.440,2.522) (3.550,2.420) (3.660,2.316) (3.770,2.210) (3.880,2.104) (3.990,1.997) (4.100,1.888) (4.210,1.780) (4.320,1.671) (4.430,1.561) (4.540,1.451) (4.650,1.341) (4.760,1.231) (4.870,1.121) (4.980,1.010) (5.090,0.900) (5.200,0.789) (5.310,0.678) (5.420,0.567) (5.530,0.457) (5.640,0.346) (5.750,0.250) (5.860,0.250) (5.970,0.250) (6.080,0.250) (6.190,0.250) (6.300,0.250)};
\node[lab,text=khmag] at (2.2,0.7) {sand};\node[lab,text=khcyan!70!black] at (5.5,3.3) {silt loam};
\end{kfig}

\section{Derivation of the closed form}
Insert the inverse retention curve, Equation~(\ref{eq:vginv}), and substitute $y=x^{1/m}$, so that $x=y^m$ and $dx=m\,y^{m-1}dy$:
\begin{equation}\label{eq:int1}\boxed{\int_0^{\Se}\frac{dx}{h(x)}=\alpha m\int_0^{\Se^{1/m}} y^{\,m-1}\bigl(y^{-1}-1\bigr)^{-1/n}dy
=\alpha m\int_0^{\Se^{1/m}} y^{\,m-1+1/n}(1-y)^{-1/n}dy .}\end{equation}
With the restriction $m=1-1/n$ the exponent of $y$ is $m-1+1/n=0$, and $-1/n=m-1$, so the integral is elementary:
\begin{equation}\label{eq:int2}\boxed{\int_0^{\Se}\frac{dx}{h(x)}=\alpha m\int_0^{\Se^{1/m}}(1-y)^{m-1}dy=\alpha\Bigl[1-\bigl(1-\Se^{1/m}\bigr)^{m}\Bigr].}\end{equation}
\par\noindent where $x$ and $y$ are dummy variables of integration, $m$ and $\alpha$ the van Genuchten parameters, and $\Se$ the effective saturation; the substitution $y=x^{1/m}$ turns the integral into the form of the incomplete beta function.

Dividing by the same integral taken to $\Se=1$, which equals $\alpha$, and inserting in Equation~(\ref{eq:mualem}), gives the van Genuchten--Mualem pair:
\begin{equation}\label{eq:vgm}\boxed{\boxed{\;\Kr(\Se)=\Se^{\ell}\Bigl[1-\bigl(1-\Se^{1/m}\bigr)^{m}\Bigr]^{2}\;}}\end{equation}
with $\ell=0.5$, the value Mualem found as the average of many soils.

\exbox{Conductivity of sand and silty clay loam}{For the sand S03 at $h=100$ cm we found $\Se=0.0112$, with $m=0.627$. Then $\Se^{1/m}=0.0112^{1.595}=7.8\times10^{-4}$, and $1-(1-7.8\times10^{-4})^{0.627}=4.9\times10^{-4}$. So $\Kr=0.0112^{0.5}\,(4.9\times10^{-4})^{2}=0.106\times2.4\times10^{-7}=2.5\times10^{-8}$. At 300 cm it is $2.7\times10^{-11}$, and at 1000 cm $1.6\times10^{-14}$. For the silty clay loam S10 at 300 cm, $\Kr=1.5\times10^{-3}$, and at 1000 cm $8.6\times10^{-5}$. At 300 cm the conductivity of the sand is $10^{8}$ times smaller than that of the silty clay loam, in relative terms. This is the physical reason why a 300 cm trigger is unreachable in sand: the crop cannot pull water through a conductivity that has collapsed.}



\fig{p4_reach}{0.9}{The reach of a trigger. (a) Irrigation against trigger suction for the ten soils. (b) The same runs on the reach scale $h_t/h_c$: the curves fall to zero near 1. (c) The closed form, Equation~(\ref{eq:reach}), against the observed reach interval from Daisy. (d) Yield: a trigger beyond the reach is rainfed.}

\begin{daisybox}{what Daisy does with the conductivity}
Daisy solves the Richards equation with these functions, so the model's conductivity at each suction follows Equation~(\ref{eq:vgm}), with its own saturated conductivity $\Ks$ per soil. Eight of the ten soils have literature values of $\Ks$ (Carsel and Parrish, 1988); the reach law of Chapter~\ref{ch:reach} will use only $\Kr$, that is, it does not need $\Ks$.
\end{daisybox}

\exbox{}{Verify Equation~(\ref{eq:int2}) by differentiating the right-hand side with respect to $\Se$.}\sol  Differentiate: $\frac{d}{d\Se}\alpha[1-(1-\Se^{1/m})^m]=\alpha(1-y)^{m-1}\Se^{1/m-1}$ with $y=\Se^{1/m}$. The integrand is $1/h=\alpha(\Se^{-1/m}-1)^{-1/n}=\alpha\,y^{1/n}(1-y)^{-1/n}$. Since $-1/n=m-1$ and $\Se^{1/m-1}=\Se^{1/(n-1)}=y^{1/n}$ (because $mn=n-1$), the two agree.

\exbox{}{Compute $\Kr$ for the loam S06 at $h=300$ cm.}\sol  $\Kr(S06,300\text{ cm})=3.8\times10^{-5}$.

\exbox{}{Show that $\Kr\to1$ as $\Se\to1$ and $\Kr\to0$ as $\Se\to0$.}\sol  At $\Se=1$ the bracket is $1-0=1$; at $\Se\to0$, $\Se^\ell\to0$ and the bracket $\to0$.

\exbox{}{Why does the restriction $m=1-1/n$ make the integral elementary? What would be needed otherwise?}\sol  The exponent $m-1+1/n$ vanishes, leaving $(1-y)^{m-1}$ whose primitive is elementary; otherwise one needs an incomplete beta function.
\IfFileExists{extra/app01.tex}{\subsection{The output of the application, panel by panel}
The water of Chapter 1 is not only a set of equations: in SWAFL it is a set of pictures that change with the settings of the user. The layer \emph{Soil hydraulic drivers} shows how the retention curve and the conductivity of the soil move when the tillage, the consolidation and the texture change, and the figures below are its panels for the coarse loam.

The first panel ranks the ten soils by the change of a hydraulic property (the saturated conductivity or the water held) when one driver is moved from its present value. The bars are the sensitivity of the property to the driver; the scatter on the right places the soils on the plane of the two drivers. It is a direct reading of the van Genuchten parameters of this chapter as functions of the pore structure.

\appfig{v7_0}{0.96}{What drives change in soil hydraulic properties: the response of the soil to tillage, consolidation rate and organic matter, soil by soil.}

The graphs follow the drainage against the consolidation rate and the water content against the suction for conventional, reduced and no-till systems; clicking a point applies that setting to the whole program, so the effect of a change on the retention curve can be seen at once.

\appfig{v7_1}{0.96}{Graphs of the drivers at your settings: drainage and water content against the consolidation rate, for the three tillage systems.}

The pictures are sections through the topsoil of the ensemble runs after 0, 50, 100 and 200 days of settling. The dark cracks and macropores close as the soil consolidates, and the pore index printed under each section (0.92 and 0.67 for two states) is the number that enters the conductivity of the layer.

\appfig{v7_2}{0.96}{Real soil profiles: how tillage and settling change the soil, layer by layer, in the pictures of the model.}

}{}
\IfFileExists{extra/long01.tex}{\subsection{Complete characterisation of the loam}
\exbox{Retention, conductivity and storage of a loam}{Characterise the loam S06 completely: find $S_e$, $\theta$ and $K_r$ at six suctions, then the water held between field capacity and wilting, and the total and readily available water of a root zone 0.8 m deep.}\sol
\ans Van Genuchten parameters of the loam: $\alpha=0.036$ cm$^{-1}$, $n=1.56$, so $m=1-1/n=0.359$; $\theta_r=0.05$, $\theta_s=0.43$
\ans At $h=10$ cm: $S_e=[1+(\alpha h)^n]^{-m}=0.9358$; $\theta=\theta_r+(\theta_s-\theta_r)S_e=0.4056$; $K_r=0.215$
\ans At $h=33$ cm: $S_e=[1+(\alpha h)^n]^{-m}=0.7406$; $\theta=\theta_r+(\theta_s-\theta_r)S_e=0.3314$; $K_r=0.0293$
\ans At $h=100$ cm: $S_e=[1+(\alpha h)^n]^{-m}=0.4663$; $\theta=\theta_r+(\theta_s-\theta_r)S_e=0.2272$; $K_r=0.00136$
\ans At $h=330$ cm: $S_e=[1+(\alpha h)^n]^{-m}=0.2482$; $\theta=\theta_r+(\theta_s-\theta_r)S_e=0.1443$; $K_r=2.77e-05$
\ans At $h=1000$ cm: $S_e=[1+(\alpha h)^n]^{-m}=0.1342$; $\theta=\theta_r+(\theta_s-\theta_r)S_e=0.1010$; $K_r=6.55e-07$
\ans At $h=15849$ cm: $S_e=[1+(\alpha h)^n]^{-m}=0.0286$; $\theta=\theta_r+(\theta_s-\theta_r)S_e=0.0609$; $K_r=5.48e-11$
\ans Field capacity (33 kPa) $\theta_{fc}=0.1443$ and wilting point (1500 kPa) $\theta_{wp}=0.0609$, so the available water is $\theta_{fc}-\theta_{wp}=0.0835$ m$^3$\,m$^{-3}$
\ans For a root zone of 0.8 m: total available water $TAW=1000(\theta_{fc}-\theta_{wp})Z_r=67$ mm
\ans Readily available water with depletion fraction $p=0.5$: $RAW=pTAW=33$ mm; at that point the suction is the trigger of Chapter~4 \finans


}{}
\IfFileExists{extra/proj01.tex}{\banner{Design project 1: the soil-water data of a field for a schedule}
\textbf{Brief.} A grower of maize on a 24-ha field of loam asks for the soil-water numbers on which an irrigation schedule is built. The roots reach 0.9 m. The profile was sampled in four layers; the retention curve of each layer was taken from the van Genuchten parameters measured on cores. The peak crop evapotranspiration is $ET_c=6.2$ mm\,d$^{-1}$, the management allowed depletion is $p=0.55$, the system is sprinklers of 75 per cent application efficiency, the pump delivers $Q=48$ L\,s$^{-1}$ and the system irrigates one tenth of the field at a time. The task is to find the total and readily available water, the net and gross irrigation depths, the interval, the volume per irrigation, the set time and the number of sets per cycle, and to check the cycle against the peak demand.

\begin{table}[!htbp]\centering\small\caption{Step 1: field capacity, wilting point and total available water of each layer, computed from the van Genuchten parameters.}
\begin{tabular}{lrrrrrrr}\toprule
\hd{Layer (cm)} & \hd{$z$ (cm)} & \hd{$\theta_s$} & \hd{$\alpha$} & \hd{$n$} & \hd{$\theta_{FC}$} & \hd{$\theta_{WP}$} & \hd{TAW (mm)}\\\midrule
0--20 & 20 & 0.468 & 0.041 & 1.52 & 0.227 & 0.035 & 38.3\\
20--40 & 20 & 0.455 & 0.036 & 1.56 & 0.223 & 0.032 & 38.1\\
40--60 & 20 & 0.447 & 0.030 & 1.52 & 0.248 & 0.039 & 41.8\\
60--90 & 30 & 0.440 & 0.026 & 1.48 & 0.269 & 0.048 & 66.4\\\bottomrule\end{tabular}\end{table}

\textbf{Step 2: allowed depletion and net depth.} The sum of the table is $\mathrm{TAW}=185$ mm. With $p=0.55$, $\mathrm{RAW}=0.55\times185=102$ mm, which is the net depth $Z_n$ applied at each irrigation.

\textbf{Step 3: interval.} At the peak demand the root zone loses $6.2$ mm\,d$^{-1}$ and the allowed depletion is used in $T=Z_n/ET_c=102/6.2=16.4$ days. The system must therefore return to each set at least every 16 days; the cycle is rounded down to 16 days.

\textbf{Step 4: gross depth and volume.} With an efficiency of 75 per cent the gross depth is $Z_g=Z_n/E_a=102/0.75=135$ mm. One set covers 2.4 ha, so the volume per set is $V=Z_g\times A=135\times10^{-3}\times24\,000=3250$ m$^3$ (1 mm over 1 ha is 10 m$^3$).

\textbf{Step 5: set time and capacity check.} At $Q=48$ L\,s$^{-1}=172.8$ m$^3$\,h$^{-1}$ the set time is $t=V/Q=3250/172.8=18.81$ h. Ten sets need $188.1$ h of pumping; with 22 h of pumping available per day, a cycle of 16 days is feasible: the available pumping time per cycle is $22\times16=352$ h against the 188 h needed.

\textbf{Conclusion.} The soil stores 102 mm of readily available water; a 75 per cent efficient system must apply 135 mm gross every 16 days, which the pump can do with a margin of 47 per cent of its pumping time.

}{}
\IfFileExists{extra/ten01.tex}{\subsection{Worked solutions on the ten soils}
The ten soils of the ensemble differ widely in their retention. This section works one question on each of them, so that the reader can see how the same formulas give different answers.

\exbox{Ten-soil example 1a: Coarse loam (S01)}{Coarse loam (S01): $\alpha=0.0961$ cm$^{-1}$, $n=1.353$, AWC $=1.5$ mm\,cm$^{-1}$. Find $m$, $S_e$ at field capacity ($h=100$ cm) and wilting point ($h=15\,849$ cm), the suction at half depletion, and TAW, RAW ($p=0.5$) and the days of supply at $ET_c=5$ mm\,d$^{-1}$ for a root zone of 0.6 m.}\sol $m=1-1/1.353=0.261$ \ans $S_e(100)=[1+(0.0961\times100)^{1.353}]^{-0.261}=0.445$; $S_e(15\,849)=0.0752$ \ans Half depletion is $f=0.5$, found by bisection of Eq.~(\ref{eq:fdep}): $h_{0.5}=471$ cm \ans $TAW=1.5\times600/10\ldots=90$ mm, i.e.\ $1.50\times60$ cm; $RAW=0.5\times90=45$ mm \ans Supply: $RAW/ET_c=45/5=9.0$ days \ans

\exbox{Ten-soil example 1a: Coarse sand (S02)}{Coarse sand (S02): $\alpha=0.0536$ cm$^{-1}$, $n=1.600$, AWC $=0.5$ mm\,cm$^{-1}$. Find $m$, $S_e$ at field capacity ($h=100$ cm) and wilting point ($h=15\,849$ cm), the suction at half depletion, and TAW, RAW ($p=0.5$) and the days of supply at $ET_c=5$ mm\,d$^{-1}$ for a root zone of 0.6 m.}\sol $m=1-1/1.600=0.375$ \ans $S_e(100)=[1+(0.0536\times100)^{1.600}]^{-0.375}=0.356$; $S_e(15\,849)=0.0175$ \ans Half depletion is $f=0.5$, found by bisection of Eq.~(\ref{eq:fdep}): $h_{0.5}=303$ cm \ans $TAW=0.5\times600/10\ldots=30$ mm, i.e.\ $0.50\times60$ cm; $RAW=0.5\times30=15$ mm \ans Supply: $RAW/ET_c=15/5=3.0$ days \ans

\exbox{Ten-soil example 1a: Sand (S03)}{Sand (S03): $\alpha=0.1450$ cm$^{-1}$, $n=2.680$, AWC $=0.6$ mm\,cm$^{-1}$. Find $m$, $S_e$ at field capacity ($h=100$ cm) and wilting point ($h=15\,849$ cm), the suction at half depletion, and TAW, RAW ($p=0.5$) and the days of supply at $ET_c=5$ mm\,d$^{-1}$ for a root zone of 0.6 m.}\sol $m=1-1/2.680=0.627$ \ans $S_e(100)=[1+(0.1450\times100)^{2.680}]^{-0.627}=0.011$; $S_e(15\,849)=0.0000$ \ans Half depletion is $f=0.5$, found by bisection of Eq.~(\ref{eq:fdep}): $h_{0.5}=151$ cm \ans $TAW=0.6\times600/10\ldots=36$ mm, i.e.\ $0.60\times60$ cm; $RAW=0.5\times36=18$ mm \ans Supply: $RAW/ET_c=18/5=3.6$ days \ans

\exbox{Ten-soil example 1a: Loamy sand (S04)}{Loamy sand (S04): $\alpha=0.1240$ cm$^{-1}$, $n=2.280$, AWC $=0.9$ mm\,cm$^{-1}$. Find $m$, $S_e$ at field capacity ($h=100$ cm) and wilting point ($h=15\,849$ cm), the suction at half depletion, and TAW, RAW ($p=0.5$) and the days of supply at $ET_c=5$ mm\,d$^{-1}$ for a root zone of 0.6 m.}\sol $m=1-1/2.280=0.561$ \ans $S_e(100)=[1+(0.1240\times100)^{2.280}]^{-0.561}=0.040$; $S_e(15\,849)=0.0001$ \ans Half depletion is $f=0.5$, found by bisection of Eq.~(\ref{eq:fdep}): $h_{0.5}=172$ cm \ans $TAW=0.9\times600/10\ldots=54$ mm, i.e.\ $0.90\times60$ cm; $RAW=0.5\times54=27$ mm \ans Supply: $RAW/ET_c=27/5=5.4$ days \ans

\exbox{Ten-soil example 1a: Sandy loam (S05)}{Sandy loam (S05): $\alpha=0.0750$ cm$^{-1}$, $n=1.890$, AWC $=1.2$ mm\,cm$^{-1}$. Find $m$, $S_e$ at field capacity ($h=100$ cm) and wilting point ($h=15\,849$ cm), the suction at half depletion, and TAW, RAW ($p=0.5$) and the days of supply at $ET_c=5$ mm\,d$^{-1}$ for a root zone of 0.6 m.}\sol $m=1-1/1.890=0.471$ \ans $S_e(100)=[1+(0.0750\times100)^{1.890}]^{-0.471}=0.165$; $S_e(15\,849)=0.0018$ \ans Half depletion is $f=0.5$, found by bisection of Eq.~(\ref{eq:fdep}): $h_{0.5}=217$ cm \ans $TAW=1.2\times600/10\ldots=72$ mm, i.e.\ $1.20\times60$ cm; $RAW=0.5\times72=36$ mm \ans Supply: $RAW/ET_c=36/5=7.2$ days \ans

\exbox{Ten-soil example 1a: Loam (S06)}{Loam (S06): $\alpha=0.0360$ cm$^{-1}$, $n=1.560$, AWC $=1.7$ mm\,cm$^{-1}$. Find $m$, $S_e$ at field capacity ($h=100$ cm) and wilting point ($h=15\,849$ cm), the suction at half depletion, and TAW, RAW ($p=0.5$) and the days of supply at $ET_c=5$ mm\,d$^{-1}$ for a root zone of 0.6 m.}\sol $m=1-1/1.560=0.359$ \ans $S_e(100)=[1+(0.0360\times100)^{1.560}]^{-0.359}=0.466$; $S_e(15\,849)=0.0286$ \ans Half depletion is $f=0.5$, found by bisection of Eq.~(\ref{eq:fdep}): $h_{0.5}=332$ cm \ans $TAW=1.7\times600/10\ldots=102$ mm, i.e.\ $1.70\times60$ cm; $RAW=0.5\times102=51$ mm \ans Supply: $RAW/ET_c=51/5=10.2$ days \ans

\exbox{Ten-soil example 1a: Silt loam (S07)}{Silt loam (S07): $\alpha=0.0200$ cm$^{-1}$, $n=1.410$, AWC $=2.0$ mm\,cm$^{-1}$. Find $m$, $S_e$ at field capacity ($h=100$ cm) and wilting point ($h=15\,849$ cm), the suction at half depletion, and TAW, RAW ($p=0.5$) and the days of supply at $ET_c=5$ mm\,d$^{-1}$ for a root zone of 0.6 m.}\sol $m=1-1/1.410=0.291$ \ans $S_e(100)=[1+(0.0200\times100)^{1.410}]^{-0.291}=0.686$; $S_e(15\,849)=0.0943$ \ans Half depletion is $f=0.5$, found by bisection of Eq.~(\ref{eq:fdep}): $h_{0.5}=483$ cm \ans $TAW=2.0\times600/10\ldots=120$ mm, i.e.\ $2.00\times60$ cm; $RAW=0.5\times120=60$ mm \ans Supply: $RAW/ET_c=60/5=12.0$ days \ans

\exbox{Ten-soil example 1a: Sandy clay loam (S08)}{Sandy clay loam (S08): $\alpha=0.0590$ cm$^{-1}$, $n=1.480$, AWC $=1.4$ mm\,cm$^{-1}$. Find $m$, $S_e$ at field capacity ($h=100$ cm) and wilting point ($h=15\,849$ cm), the suction at half depletion, and TAW, RAW ($p=0.5$) and the days of supply at $ET_c=5$ mm\,d$^{-1}$ for a root zone of 0.6 m.}\sol $m=1-1/1.480=0.324$ \ans $S_e(100)=[1+(0.0590\times100)^{1.480}]^{-0.324}=0.417$; $S_e(15\,849)=0.0375$ \ans Half depletion is $f=0.5$, found by bisection of Eq.~(\ref{eq:fdep}): $h_{0.5}=369$ cm \ans $TAW=1.4\times600/10\ldots=84$ mm, i.e.\ $1.40\times60$ cm; $RAW=0.5\times84=42$ mm \ans Supply: $RAW/ET_c=42/5=8.4$ days \ans

\exbox{Ten-soil example 1a: Clay loam (S09)}{Clay loam (S09): $\alpha=0.0190$ cm$^{-1}$, $n=1.310$, AWC $=1.6$ mm\,cm$^{-1}$. Find $m$, $S_e$ at field capacity ($h=100$ cm) and wilting point ($h=15\,849$ cm), the suction at half depletion, and TAW, RAW ($p=0.5$) and the days of supply at $ET_c=5$ mm\,d$^{-1}$ for a root zone of 0.6 m.}\sol $m=1-1/1.310=0.237$ \ans $S_e(100)=[1+(0.0190\times100)^{1.310}]^{-0.237}=0.753$; $S_e(15\,849)=0.1704$ \ans Half depletion is $f=0.5$, found by bisection of Eq.~(\ref{eq:fdep}): $h_{0.5}=618$ cm \ans $TAW=1.6\times600/10\ldots=96$ mm, i.e.\ $1.60\times60$ cm; $RAW=0.5\times96=48$ mm \ans Supply: $RAW/ET_c=48/5=9.6$ days \ans

\exbox{Ten-soil example 1a: Silty clay loam (S10)}{Silty clay loam (S10): $\alpha=0.0100$ cm$^{-1}$, $n=1.230$, AWC $=1.7$ mm\,cm$^{-1}$. Find $m$, $S_e$ at field capacity ($h=100$ cm) and wilting point ($h=15\,849$ cm), the suction at half depletion, and TAW, RAW ($p=0.5$) and the days of supply at $ET_c=5$ mm\,d$^{-1}$ for a root zone of 0.6 m.}\sol $m=1-1/1.230=0.187$ \ans $S_e(100)=[1+(0.0100\times100)^{1.230}]^{-0.187}=0.878$; $S_e(15\,849)=0.3118$ \ans Half depletion is $f=0.5$, found by bisection of Eq.~(\ref{eq:fdep}): $h_{0.5}=906$ cm \ans $TAW=1.7\times600/10\ldots=102$ mm, i.e.\ $1.70\times60$ cm; $RAW=0.5\times102=51$ mm \ans Supply: $RAW/ET_c=51/5=10.2$ days \ans

}{}
\IfFileExists{extra/tenb01.tex}{\subsection{Worked solutions on ten crops}
The ten crops of the library differ in their demand, their roots and their sensitivity. The same type of question is answered for each of them at $ET_0=5.5$ mm\,d$^{-1}$; the data are typical values of the crops (Doorenbos--Kassam and FAO 56).

\exbox{Crop example 1a: Maize}{Maize grows on the loam S06 (AWC 1.7 mm\,cm$^{-1}$) with a root zone of 100 cm and a depletion fraction $p=0.55$. Find TAW, RAW and the days between irrigations at its mid-season demand $K_{c}=1.20$, $ET_0=5.5$ mm\,d$^{-1}$.}\sol \givenline{$p=0.55$ ; $K_{c}=1.20$ ; $ET_0=5.5$ mm\,d$^{-1}$}\ans \[TAW=\mathrm{AWC}\,z_r=1.7\times100=170\ \text{mm}\] \ans \[RAW=p\,TAW=0.55\times170=94\ \text{mm}\] \ans \[ET_c=K_cET_0=1.20\times5.5=6.60\ \text{mm\,d$^{-1}$}\] \ans Interval $=RAW/ET_c=94/6.60=14.2$ d \finans

\exbox{Crop example 1a: Wheat}{Wheat grows on the loam S06 (AWC 1.7 mm\,cm$^{-1}$) with a root zone of 120 cm and a depletion fraction $p=0.55$. Find TAW, RAW and the days between irrigations at its mid-season demand $K_{c}=1.15$, $ET_0=5.5$ mm\,d$^{-1}$.}\sol \givenline{$p=0.55$ ; $K_{c}=1.15$ ; $ET_0=5.5$ mm\,d$^{-1}$}\ans \[TAW=\mathrm{AWC}\,z_r=1.7\times120=204\ \text{mm}\] \ans \[RAW=p\,TAW=0.55\times204=112\ \text{mm}\] \ans \[ET_c=K_cET_0=1.15\times5.5=6.32\ \text{mm\,d$^{-1}$}\] \ans Interval $=RAW/ET_c=112/6.32=17.7$ d \finans

\exbox{Crop example 1a: Cotton}{Cotton grows on the loam S06 (AWC 1.7 mm\,cm$^{-1}$) with a root zone of 110 cm and a depletion fraction $p=0.65$. Find TAW, RAW and the days between irrigations at its mid-season demand $K_{c}=1.15$, $ET_0=5.5$ mm\,d$^{-1}$.}\sol \givenline{$p=0.65$ ; $K_{c}=1.15$ ; $ET_0=5.5$ mm\,d$^{-1}$}\ans \[TAW=\mathrm{AWC}\,z_r=1.7\times110=187\ \text{mm}\] \ans \[RAW=p\,TAW=0.65\times187=122\ \text{mm}\] \ans \[ET_c=K_cET_0=1.15\times5.5=6.32\ \text{mm\,d$^{-1}$}\] \ans Interval $=RAW/ET_c=122/6.32=19.2$ d \finans

\exbox{Crop example 1a: Tomato}{Tomato grows on the loam S06 (AWC 1.7 mm\,cm$^{-1}$) with a root zone of 70 cm and a depletion fraction $p=0.40$. Find TAW, RAW and the days between irrigations at its mid-season demand $K_{c}=1.15$, $ET_0=5.5$ mm\,d$^{-1}$.}\sol \givenline{$p=0.40$ ; $K_{c}=1.15$ ; $ET_0=5.5$ mm\,d$^{-1}$}\ans \[TAW=\mathrm{AWC}\,z_r=1.7\times70=119\ \text{mm}\] \ans \[RAW=p\,TAW=0.40\times119=48\ \text{mm}\] \ans \[ET_c=K_cET_0=1.15\times5.5=6.32\ \text{mm\,d$^{-1}$}\] \ans Interval $=RAW/ET_c=48/6.32=7.5$ d \finans

\exbox{Crop example 1a: Potato}{Potato grows on the loam S06 (AWC 1.7 mm\,cm$^{-1}$) with a root zone of 50 cm and a depletion fraction $p=0.35$. Find TAW, RAW and the days between irrigations at its mid-season demand $K_{c}=1.15$, $ET_0=5.5$ mm\,d$^{-1}$.}\sol \givenline{$p=0.35$ ; $K_{c}=1.15$ ; $ET_0=5.5$ mm\,d$^{-1}$}\ans \[TAW=\mathrm{AWC}\,z_r=1.7\times50=85\ \text{mm}\] \ans \[RAW=p\,TAW=0.35\times85=30\ \text{mm}\] \ans \[ET_c=K_cET_0=1.15\times5.5=6.32\ \text{mm\,d$^{-1}$}\] \ans Interval $=RAW/ET_c=30/6.32=4.7$ d \finans

\exbox{Crop example 1a: Soybean}{Soybean grows on the loam S06 (AWC 1.7 mm\,cm$^{-1}$) with a root zone of 60 cm and a depletion fraction $p=0.50$. Find TAW, RAW and the days between irrigations at its mid-season demand $K_{c}=1.15$, $ET_0=5.5$ mm\,d$^{-1}$.}\sol \givenline{$p=0.50$ ; $K_{c}=1.15$ ; $ET_0=5.5$ mm\,d$^{-1}$}\ans \[TAW=\mathrm{AWC}\,z_r=1.7\times60=102\ \text{mm}\] \ans \[RAW=p\,TAW=0.50\times102=51\ \text{mm}\] \ans \[ET_c=K_cET_0=1.15\times5.5=6.32\ \text{mm\,d$^{-1}$}\] \ans Interval $=RAW/ET_c=51/6.32=8.1$ d \finans

\exbox{Crop example 1a: Sugar beet}{Sugar beet grows on the loam S06 (AWC 1.7 mm\,cm$^{-1}$) with a root zone of 90 cm and a depletion fraction $p=0.55$. Find TAW, RAW and the days between irrigations at its mid-season demand $K_{c}=1.20$, $ET_0=5.5$ mm\,d$^{-1}$.}\sol \givenline{$p=0.55$ ; $K_{c}=1.20$ ; $ET_0=5.5$ mm\,d$^{-1}$}\ans \[TAW=\mathrm{AWC}\,z_r=1.7\times90=153\ \text{mm}\] \ans \[RAW=p\,TAW=0.55\times153=84\ \text{mm}\] \ans \[ET_c=K_cET_0=1.20\times5.5=6.60\ \text{mm\,d$^{-1}$}\] \ans Interval $=RAW/ET_c=84/6.60=12.8$ d \finans

\exbox{Crop example 1a: Sunflower}{Sunflower grows on the loam S06 (AWC 1.7 mm\,cm$^{-1}$) with a root zone of 80 cm and a depletion fraction $p=0.45$. Find TAW, RAW and the days between irrigations at its mid-season demand $K_{c}=1.10$, $ET_0=5.5$ mm\,d$^{-1}$.}\sol \givenline{$p=0.45$ ; $K_{c}=1.10$ ; $ET_0=5.5$ mm\,d$^{-1}$}\ans \[TAW=\mathrm{AWC}\,z_r=1.7\times80=136\ \text{mm}\] \ans \[RAW=p\,TAW=0.45\times136=61\ \text{mm}\] \ans \[ET_c=K_cET_0=1.10\times5.5=6.05\ \text{mm\,d$^{-1}$}\] \ans Interval $=RAW/ET_c=61/6.05=10.1$ d \finans

\exbox{Crop example 1a: Onion}{Onion grows on the loam S06 (AWC 1.7 mm\,cm$^{-1}$) with a root zone of 35 cm and a depletion fraction $p=0.30$. Find TAW, RAW and the days between irrigations at its mid-season demand $K_{c}=1.05$, $ET_0=5.5$ mm\,d$^{-1}$.}\sol \givenline{$p=0.30$ ; $K_{c}=1.05$ ; $ET_0=5.5$ mm\,d$^{-1}$}\ans \[TAW=\mathrm{AWC}\,z_r=1.7\times35=60\ \text{mm}\] \ans \[RAW=p\,TAW=0.30\times60=18\ \text{mm}\] \ans \[ET_c=K_cET_0=1.05\times5.5=5.78\ \text{mm\,d$^{-1}$}\] \ans Interval $=RAW/ET_c=18/5.78=3.1$ d \finans

\exbox{Crop example 1a: Alfalfa}{Alfalfa grows on the loam S06 (AWC 1.7 mm\,cm$^{-1}$) with a root zone of 100 cm and a depletion fraction $p=0.55$. Find TAW, RAW and the days between irrigations at its mid-season demand $K_{c}=1.15$, $ET_0=5.5$ mm\,d$^{-1}$.}\sol \givenline{$p=0.55$ ; $K_{c}=1.15$ ; $ET_0=5.5$ mm\,d$^{-1}$}\ans \[TAW=\mathrm{AWC}\,z_r=1.7\times100=170\ \text{mm}\] \ans \[RAW=p\,TAW=0.55\times170=94\ \text{mm}\] \ans \[ET_c=K_cET_0=1.15\times5.5=6.32\ \text{mm\,d$^{-1}$}\] \ans Interval $=RAW/ET_c=94/6.32=14.8$ d \finans

}{}
\IfFileExists{extra/teny01.tex}{\subsection{Worked solutions on the ten soils (third scenario)}
The same questions are repeated for different working conditions, to show how the answer moves with them.

\exbox{Ten-soil example 1c: Coarse loam (S01)}{Coarse loam (S01): $\alpha=0.0961$ cm$^{-1}$, $n=1.353$, AWC $=1.5$ mm\,cm$^{-1}$. Find $m$, $S_e$ at field capacity ($h=100$ cm) and wilting point ($h=15\,849$ cm), the suction at half depletion, and TAW, RAW ($p=0.5$) and the days of supply at $ET_c=5$ mm\,d$^{-1}$ for a root zone of 0.6 m.}\sol \givenline{$\alpha=0.0961$ cm$^{-1}$ ; $n=1.353$ ; $=1.5$ mm\,cm$^{-1}$ ; $h=100$ cm ; $h=15\,849$ cm ; $p=0.5$ ; $ET_c=5$ mm\,d$^{-1}$}\ans $m=1-1/1.353=0.261$ \ans $S_e(100)=[1+(0.0961\times100)^{1.353}]^{-0.261}=0.445$; $S_e(15\,849)=0.0752$ \ans Half depletion is $f=0.5$, found by bisection of Eq.~(\ref{eq:fdep}): $h_{0.5}=471$ cm \ans $TAW=1.5\times600/10\ldots=90$ mm, i.e.\ $1.50\times60$ cm; $RAW=0.5\times90=45$ mm \ans Supply: $RAW/ET_c=45/5=9.0$ days \finans

\exbox{Ten-soil example 1c: Coarse sand (S02)}{Coarse sand (S02): $\alpha=0.0536$ cm$^{-1}$, $n=1.600$, AWC $=0.5$ mm\,cm$^{-1}$. Find $m$, $S_e$ at field capacity ($h=100$ cm) and wilting point ($h=15\,849$ cm), the suction at half depletion, and TAW, RAW ($p=0.5$) and the days of supply at $ET_c=5$ mm\,d$^{-1}$ for a root zone of 0.6 m.}\sol \givenline{$\alpha=0.0536$ cm$^{-1}$ ; $n=1.600$ ; $=0.5$ mm\,cm$^{-1}$ ; $h=100$ cm ; $h=15\,849$ cm ; $p=0.5$ ; $ET_c=5$ mm\,d$^{-1}$}\ans $m=1-1/1.600=0.375$ \ans $S_e(100)=[1+(0.0536\times100)^{1.600}]^{-0.375}=0.356$; $S_e(15\,849)=0.0175$ \ans Half depletion is $f=0.5$, found by bisection of Eq.~(\ref{eq:fdep}): $h_{0.5}=303$ cm \ans $TAW=0.5\times600/10\ldots=30$ mm, i.e.\ $0.50\times60$ cm; $RAW=0.5\times30=15$ mm \ans Supply: $RAW/ET_c=15/5=3.0$ days \finans

\exbox{Ten-soil example 1c: Sand (S03)}{Sand (S03): $\alpha=0.1450$ cm$^{-1}$, $n=2.680$, AWC $=0.6$ mm\,cm$^{-1}$. Find $m$, $S_e$ at field capacity ($h=100$ cm) and wilting point ($h=15\,849$ cm), the suction at half depletion, and TAW, RAW ($p=0.5$) and the days of supply at $ET_c=5$ mm\,d$^{-1}$ for a root zone of 0.6 m.}\sol \givenline{$\alpha=0.1450$ cm$^{-1}$ ; $n=2.680$ ; $=0.6$ mm\,cm$^{-1}$ ; $h=100$ cm ; $h=15\,849$ cm ; $p=0.5$ ; $ET_c=5$ mm\,d$^{-1}$}\ans $m=1-1/2.680=0.627$ \ans $S_e(100)=[1+(0.1450\times100)^{2.680}]^{-0.627}=0.011$; $S_e(15\,849)=0.0000$ \ans Half depletion is $f=0.5$, found by bisection of Eq.~(\ref{eq:fdep}): $h_{0.5}=151$ cm \ans $TAW=0.6\times600/10\ldots=36$ mm, i.e.\ $0.60\times60$ cm; $RAW=0.5\times36=18$ mm \ans Supply: $RAW/ET_c=18/5=3.6$ days \finans

\exbox{Ten-soil example 1c: Loamy sand (S04)}{Loamy sand (S04): $\alpha=0.1240$ cm$^{-1}$, $n=2.280$, AWC $=0.9$ mm\,cm$^{-1}$. Find $m$, $S_e$ at field capacity ($h=100$ cm) and wilting point ($h=15\,849$ cm), the suction at half depletion, and TAW, RAW ($p=0.5$) and the days of supply at $ET_c=5$ mm\,d$^{-1}$ for a root zone of 0.6 m.}\sol \givenline{$\alpha=0.1240$ cm$^{-1}$ ; $n=2.280$ ; $=0.9$ mm\,cm$^{-1}$ ; $h=100$ cm ; $h=15\,849$ cm ; $p=0.5$ ; $ET_c=5$ mm\,d$^{-1}$}\ans $m=1-1/2.280=0.561$ \ans $S_e(100)=[1+(0.1240\times100)^{2.280}]^{-0.561}=0.040$; $S_e(15\,849)=0.0001$ \ans Half depletion is $f=0.5$, found by bisection of Eq.~(\ref{eq:fdep}): $h_{0.5}=172$ cm \ans $TAW=0.9\times600/10\ldots=54$ mm, i.e.\ $0.90\times60$ cm; $RAW=0.5\times54=27$ mm \ans Supply: $RAW/ET_c=27/5=5.4$ days \finans

\exbox{Ten-soil example 1c: Sandy loam (S05)}{Sandy loam (S05): $\alpha=0.0750$ cm$^{-1}$, $n=1.890$, AWC $=1.2$ mm\,cm$^{-1}$. Find $m$, $S_e$ at field capacity ($h=100$ cm) and wilting point ($h=15\,849$ cm), the suction at half depletion, and TAW, RAW ($p=0.5$) and the days of supply at $ET_c=5$ mm\,d$^{-1}$ for a root zone of 0.6 m.}\sol \givenline{$\alpha=0.0750$ cm$^{-1}$ ; $n=1.890$ ; $=1.2$ mm\,cm$^{-1}$ ; $h=100$ cm ; $h=15\,849$ cm ; $p=0.5$ ; $ET_c=5$ mm\,d$^{-1}$}\ans $m=1-1/1.890=0.471$ \ans $S_e(100)=[1+(0.0750\times100)^{1.890}]^{-0.471}=0.165$; $S_e(15\,849)=0.0018$ \ans Half depletion is $f=0.5$, found by bisection of Eq.~(\ref{eq:fdep}): $h_{0.5}=217$ cm \ans $TAW=1.2\times600/10\ldots=72$ mm, i.e.\ $1.20\times60$ cm; $RAW=0.5\times72=36$ mm \ans Supply: $RAW/ET_c=36/5=7.2$ days \finans

\exbox{Ten-soil example 1c: Loam (S06)}{Loam (S06): $\alpha=0.0360$ cm$^{-1}$, $n=1.560$, AWC $=1.7$ mm\,cm$^{-1}$. Find $m$, $S_e$ at field capacity ($h=100$ cm) and wilting point ($h=15\,849$ cm), the suction at half depletion, and TAW, RAW ($p=0.5$) and the days of supply at $ET_c=5$ mm\,d$^{-1}$ for a root zone of 0.6 m.}\sol \givenline{$\alpha=0.0360$ cm$^{-1}$ ; $n=1.560$ ; $=1.7$ mm\,cm$^{-1}$ ; $h=100$ cm ; $h=15\,849$ cm ; $p=0.5$ ; $ET_c=5$ mm\,d$^{-1}$}\ans $m=1-1/1.560=0.359$ \ans $S_e(100)=[1+(0.0360\times100)^{1.560}]^{-0.359}=0.466$; $S_e(15\,849)=0.0286$ \ans Half depletion is $f=0.5$, found by bisection of Eq.~(\ref{eq:fdep}): $h_{0.5}=332$ cm \ans $TAW=1.7\times600/10\ldots=102$ mm, i.e.\ $1.70\times60$ cm; $RAW=0.5\times102=51$ mm \ans Supply: $RAW/ET_c=51/5=10.2$ days \finans

\exbox{Ten-soil example 1c: Silt loam (S07)}{Silt loam (S07): $\alpha=0.0200$ cm$^{-1}$, $n=1.410$, AWC $=2.0$ mm\,cm$^{-1}$. Find $m$, $S_e$ at field capacity ($h=100$ cm) and wilting point ($h=15\,849$ cm), the suction at half depletion, and TAW, RAW ($p=0.5$) and the days of supply at $ET_c=5$ mm\,d$^{-1}$ for a root zone of 0.6 m.}\sol \givenline{$\alpha=0.0200$ cm$^{-1}$ ; $n=1.410$ ; $=2.0$ mm\,cm$^{-1}$ ; $h=100$ cm ; $h=15\,849$ cm ; $p=0.5$ ; $ET_c=5$ mm\,d$^{-1}$}\ans $m=1-1/1.410=0.291$ \ans $S_e(100)=[1+(0.0200\times100)^{1.410}]^{-0.291}=0.686$; $S_e(15\,849)=0.0943$ \ans Half depletion is $f=0.5$, found by bisection of Eq.~(\ref{eq:fdep}): $h_{0.5}=483$ cm \ans $TAW=2.0\times600/10\ldots=120$ mm, i.e.\ $2.00\times60$ cm; $RAW=0.5\times120=60$ mm \ans Supply: $RAW/ET_c=60/5=12.0$ days \finans

\exbox{Ten-soil example 1c: Sandy clay loam (S08)}{Sandy clay loam (S08): $\alpha=0.0590$ cm$^{-1}$, $n=1.480$, AWC $=1.4$ mm\,cm$^{-1}$. Find $m$, $S_e$ at field capacity ($h=100$ cm) and wilting point ($h=15\,849$ cm), the suction at half depletion, and TAW, RAW ($p=0.5$) and the days of supply at $ET_c=5$ mm\,d$^{-1}$ for a root zone of 0.6 m.}\sol \givenline{$\alpha=0.0590$ cm$^{-1}$ ; $n=1.480$ ; $=1.4$ mm\,cm$^{-1}$ ; $h=100$ cm ; $h=15\,849$ cm ; $p=0.5$ ; $ET_c=5$ mm\,d$^{-1}$}\ans $m=1-1/1.480=0.324$ \ans $S_e(100)=[1+(0.0590\times100)^{1.480}]^{-0.324}=0.417$; $S_e(15\,849)=0.0375$ \ans Half depletion is $f=0.5$, found by bisection of Eq.~(\ref{eq:fdep}): $h_{0.5}=369$ cm \ans $TAW=1.4\times600/10\ldots=84$ mm, i.e.\ $1.40\times60$ cm; $RAW=0.5\times84=42$ mm \ans Supply: $RAW/ET_c=42/5=8.4$ days \finans

\exbox{Ten-soil example 1c: Clay loam (S09)}{Clay loam (S09): $\alpha=0.0190$ cm$^{-1}$, $n=1.310$, AWC $=1.6$ mm\,cm$^{-1}$. Find $m$, $S_e$ at field capacity ($h=100$ cm) and wilting point ($h=15\,849$ cm), the suction at half depletion, and TAW, RAW ($p=0.5$) and the days of supply at $ET_c=5$ mm\,d$^{-1}$ for a root zone of 0.6 m.}\sol \givenline{$\alpha=0.0190$ cm$^{-1}$ ; $n=1.310$ ; $=1.6$ mm\,cm$^{-1}$ ; $h=100$ cm ; $h=15\,849$ cm ; $p=0.5$ ; $ET_c=5$ mm\,d$^{-1}$}\ans $m=1-1/1.310=0.237$ \ans $S_e(100)=[1+(0.0190\times100)^{1.310}]^{-0.237}=0.753$; $S_e(15\,849)=0.1704$ \ans Half depletion is $f=0.5$, found by bisection of Eq.~(\ref{eq:fdep}): $h_{0.5}=618$ cm \ans $TAW=1.6\times600/10\ldots=96$ mm, i.e.\ $1.60\times60$ cm; $RAW=0.5\times96=48$ mm \ans Supply: $RAW/ET_c=48/5=9.6$ days \finans

\exbox{Ten-soil example 1c: Silty clay loam (S10)}{Silty clay loam (S10): $\alpha=0.0100$ cm$^{-1}$, $n=1.230$, AWC $=1.7$ mm\,cm$^{-1}$. Find $m$, $S_e$ at field capacity ($h=100$ cm) and wilting point ($h=15\,849$ cm), the suction at half depletion, and TAW, RAW ($p=0.5$) and the days of supply at $ET_c=5$ mm\,d$^{-1}$ for a root zone of 0.6 m.}\sol \givenline{$\alpha=0.0100$ cm$^{-1}$ ; $n=1.230$ ; $=1.7$ mm\,cm$^{-1}$ ; $h=100$ cm ; $h=15\,849$ cm ; $p=0.5$ ; $ET_c=5$ mm\,d$^{-1}$}\ans $m=1-1/1.230=0.187$ \ans $S_e(100)=[1+(0.0100\times100)^{1.230}]^{-0.187}=0.878$; $S_e(15\,849)=0.3118$ \ans Half depletion is $f=0.5$, found by bisection of Eq.~(\ref{eq:fdep}): $h_{0.5}=906$ cm \ans $TAW=1.7\times600/10\ldots=102$ mm, i.e.\ $1.70\times60$ cm; $RAW=0.5\times102=51$ mm \ans Supply: $RAW/ET_c=51/5=10.2$ days \finans

}{}
\IfFileExists{extra/tenby01.tex}{\subsection{Worked solutions on ten crops (hot climate)}
The ten crops of the library differ in their demand, their roots and their sensitivity. The same type of question is answered for each of them at $ET_0=7.5$ mm\,d$^{-1}$; the data are typical values of the crops (Doorenbos--Kassam and FAO 56).

\exbox{Crop example 1c: Maize}{Maize grows on the loam S06 (AWC 1.7 mm\,cm$^{-1}$) with a root zone of 100 cm and a depletion fraction $p=0.55$. Find TAW, RAW and the days between irrigations at its mid-season demand $K_{c}=1.20$, $ET_0=7.5$ mm\,d$^{-1}$.}\sol \givenline{$p=0.55$ ; $K_{c}=1.20$ ; $ET_0=7.5$ mm\,d$^{-1}$}\ans \[TAW=\mathrm{AWC}\,z_r=1.7\times100=170\ \text{mm}\] \ans \[RAW=p\,TAW=0.55\times170=94\ \text{mm}\] \ans \[ET_c=K_cET_0=1.20\times7.5=9.00\ \text{mm\,d$^{-1}$}\] \ans Interval $=RAW/ET_c=94/9.00=10.4$ d \finans

\exbox{Crop example 1c: Wheat}{Wheat grows on the loam S06 (AWC 1.7 mm\,cm$^{-1}$) with a root zone of 120 cm and a depletion fraction $p=0.55$. Find TAW, RAW and the days between irrigations at its mid-season demand $K_{c}=1.15$, $ET_0=7.5$ mm\,d$^{-1}$.}\sol \givenline{$p=0.55$ ; $K_{c}=1.15$ ; $ET_0=7.5$ mm\,d$^{-1}$}\ans \[TAW=\mathrm{AWC}\,z_r=1.7\times120=204\ \text{mm}\] \ans \[RAW=p\,TAW=0.55\times204=112\ \text{mm}\] \ans \[ET_c=K_cET_0=1.15\times7.5=8.62\ \text{mm\,d$^{-1}$}\] \ans Interval $=RAW/ET_c=112/8.62=13.0$ d \finans

\exbox{Crop example 1c: Cotton}{Cotton grows on the loam S06 (AWC 1.7 mm\,cm$^{-1}$) with a root zone of 110 cm and a depletion fraction $p=0.65$. Find TAW, RAW and the days between irrigations at its mid-season demand $K_{c}=1.15$, $ET_0=7.5$ mm\,d$^{-1}$.}\sol \givenline{$p=0.65$ ; $K_{c}=1.15$ ; $ET_0=7.5$ mm\,d$^{-1}$}\ans \[TAW=\mathrm{AWC}\,z_r=1.7\times110=187\ \text{mm}\] \ans \[RAW=p\,TAW=0.65\times187=122\ \text{mm}\] \ans \[ET_c=K_cET_0=1.15\times7.5=8.62\ \text{mm\,d$^{-1}$}\] \ans Interval $=RAW/ET_c=122/8.62=14.1$ d \finans

\exbox{Crop example 1c: Tomato}{Tomato grows on the loam S06 (AWC 1.7 mm\,cm$^{-1}$) with a root zone of 70 cm and a depletion fraction $p=0.40$. Find TAW, RAW and the days between irrigations at its mid-season demand $K_{c}=1.15$, $ET_0=7.5$ mm\,d$^{-1}$.}\sol \givenline{$p=0.40$ ; $K_{c}=1.15$ ; $ET_0=7.5$ mm\,d$^{-1}$}\ans \[TAW=\mathrm{AWC}\,z_r=1.7\times70=119\ \text{mm}\] \ans \[RAW=p\,TAW=0.40\times119=48\ \text{mm}\] \ans \[ET_c=K_cET_0=1.15\times7.5=8.62\ \text{mm\,d$^{-1}$}\] \ans Interval $=RAW/ET_c=48/8.62=5.5$ d \finans

\exbox{Crop example 1c: Potato}{Potato grows on the loam S06 (AWC 1.7 mm\,cm$^{-1}$) with a root zone of 50 cm and a depletion fraction $p=0.35$. Find TAW, RAW and the days between irrigations at its mid-season demand $K_{c}=1.15$, $ET_0=7.5$ mm\,d$^{-1}$.}\sol \givenline{$p=0.35$ ; $K_{c}=1.15$ ; $ET_0=7.5$ mm\,d$^{-1}$}\ans \[TAW=\mathrm{AWC}\,z_r=1.7\times50=85\ \text{mm}\] \ans \[RAW=p\,TAW=0.35\times85=30\ \text{mm}\] \ans \[ET_c=K_cET_0=1.15\times7.5=8.62\ \text{mm\,d$^{-1}$}\] \ans Interval $=RAW/ET_c=30/8.62=3.4$ d \finans

\exbox{Crop example 1c: Soybean}{Soybean grows on the loam S06 (AWC 1.7 mm\,cm$^{-1}$) with a root zone of 60 cm and a depletion fraction $p=0.50$. Find TAW, RAW and the days between irrigations at its mid-season demand $K_{c}=1.15$, $ET_0=7.5$ mm\,d$^{-1}$.}\sol \givenline{$p=0.50$ ; $K_{c}=1.15$ ; $ET_0=7.5$ mm\,d$^{-1}$}\ans \[TAW=\mathrm{AWC}\,z_r=1.7\times60=102\ \text{mm}\] \ans \[RAW=p\,TAW=0.50\times102=51\ \text{mm}\] \ans \[ET_c=K_cET_0=1.15\times7.5=8.62\ \text{mm\,d$^{-1}$}\] \ans Interval $=RAW/ET_c=51/8.62=5.9$ d \finans

\exbox{Crop example 1c: Sugar beet}{Sugar beet grows on the loam S06 (AWC 1.7 mm\,cm$^{-1}$) with a root zone of 90 cm and a depletion fraction $p=0.55$. Find TAW, RAW and the days between irrigations at its mid-season demand $K_{c}=1.20$, $ET_0=7.5$ mm\,d$^{-1}$.}\sol \givenline{$p=0.55$ ; $K_{c}=1.20$ ; $ET_0=7.5$ mm\,d$^{-1}$}\ans \[TAW=\mathrm{AWC}\,z_r=1.7\times90=153\ \text{mm}\] \ans \[RAW=p\,TAW=0.55\times153=84\ \text{mm}\] \ans \[ET_c=K_cET_0=1.20\times7.5=9.00\ \text{mm\,d$^{-1}$}\] \ans Interval $=RAW/ET_c=84/9.00=9.4$ d \finans

\exbox{Crop example 1c: Sunflower}{Sunflower grows on the loam S06 (AWC 1.7 mm\,cm$^{-1}$) with a root zone of 80 cm and a depletion fraction $p=0.45$. Find TAW, RAW and the days between irrigations at its mid-season demand $K_{c}=1.10$, $ET_0=7.5$ mm\,d$^{-1}$.}\sol \givenline{$p=0.45$ ; $K_{c}=1.10$ ; $ET_0=7.5$ mm\,d$^{-1}$}\ans \[TAW=\mathrm{AWC}\,z_r=1.7\times80=136\ \text{mm}\] \ans \[RAW=p\,TAW=0.45\times136=61\ \text{mm}\] \ans \[ET_c=K_cET_0=1.10\times7.5=8.25\ \text{mm\,d$^{-1}$}\] \ans Interval $=RAW/ET_c=61/8.25=7.4$ d \finans

\exbox{Crop example 1c: Onion}{Onion grows on the loam S06 (AWC 1.7 mm\,cm$^{-1}$) with a root zone of 35 cm and a depletion fraction $p=0.30$. Find TAW, RAW and the days between irrigations at its mid-season demand $K_{c}=1.05$, $ET_0=7.5$ mm\,d$^{-1}$.}\sol \givenline{$p=0.30$ ; $K_{c}=1.05$ ; $ET_0=7.5$ mm\,d$^{-1}$}\ans \[TAW=\mathrm{AWC}\,z_r=1.7\times35=60\ \text{mm}\] \ans \[RAW=p\,TAW=0.30\times60=18\ \text{mm}\] \ans \[ET_c=K_cET_0=1.05\times7.5=7.88\ \text{mm\,d$^{-1}$}\] \ans Interval $=RAW/ET_c=18/7.88=2.3$ d \finans

\exbox{Crop example 1c: Alfalfa}{Alfalfa grows on the loam S06 (AWC 1.7 mm\,cm$^{-1}$) with a root zone of 100 cm and a depletion fraction $p=0.55$. Find TAW, RAW and the days between irrigations at its mid-season demand $K_{c}=1.15$, $ET_0=7.5$ mm\,d$^{-1}$.}\sol \givenline{$p=0.55$ ; $K_{c}=1.15$ ; $ET_0=7.5$ mm\,d$^{-1}$}\ans \[TAW=\mathrm{AWC}\,z_r=1.7\times100=170\ \text{mm}\] \ans \[RAW=p\,TAW=0.55\times170=94\ \text{mm}\] \ans \[ET_c=K_cET_0=1.15\times7.5=8.62\ \text{mm\,d$^{-1}$}\] \ans Interval $=RAW/ET_c=94/8.62=10.8$ d \finans

}{}
\IfFileExists{extra/tenx01.tex}{\subsection{Worked solutions on the ten soils (second scenario)}
The same questions are repeated for different working conditions, to show how the answer moves with them.

\exbox{Ten-soil example 1b: Coarse loam (S01)}{Coarse loam (S01): $\alpha=0.0961$ cm$^{-1}$, $n=1.353$, AWC $=1.5$ mm\,cm$^{-1}$. Find $m$, $S_e$ at field capacity ($h=100$ cm) and wilting point ($h=15\,849$ cm), the suction at half depletion, and TAW, RAW ($p=0.5$) and the days of supply at $ET_c=5$ mm\,d$^{-1}$ for a root zone of 0.6 m.}\sol \givenline{$\alpha=0.0961$ cm$^{-1}$ ; $n=1.353$ ; $=1.5$ mm\,cm$^{-1}$ ; $h=100$ cm ; $h=15\,849$ cm ; $p=0.5$ ; $ET_c=5$ mm\,d$^{-1}$}\ans $m=1-1/1.353=0.261$ \ans \[S_e(100)=[1+(0.0961\times100)^{1.353}]^{-0.261}=0.445\]
$S_e(15\,849)=0.0752$ \ans Half depletion is $f=0.5$, found by bisection of Eq.~(\ref{eq:fdep}): $h_{0.5}=471$ cm \ans \[TAW=1.5\times600/10\ldots=90\ \text{mm}\]
i.e.\ $1.50\times60$ cm; $RAW=0.5\times90=45$ mm \ans Supply: $RAW/ET_c=45/5=9.0$ days \finans

\exbox{Ten-soil example 1b: Coarse sand (S02)}{Coarse sand (S02): $\alpha=0.0536$ cm$^{-1}$, $n=1.600$, AWC $=0.5$ mm\,cm$^{-1}$. Find $m$, $S_e$ at field capacity ($h=100$ cm) and wilting point ($h=15\,849$ cm), the suction at half depletion, and TAW, RAW ($p=0.5$) and the days of supply at $ET_c=5$ mm\,d$^{-1}$ for a root zone of 0.6 m.}\sol \givenline{$\alpha=0.0536$ cm$^{-1}$ ; $n=1.600$ ; $=0.5$ mm\,cm$^{-1}$ ; $h=100$ cm ; $h=15\,849$ cm ; $p=0.5$ ; $ET_c=5$ mm\,d$^{-1}$}\ans $m=1-1/1.600=0.375$ \ans \[S_e(100)=[1+(0.0536\times100)^{1.600}]^{-0.375}=0.356\]
$S_e(15\,849)=0.0175$ \ans Half depletion is $f=0.5$, found by bisection of Eq.~(\ref{eq:fdep}): $h_{0.5}=303$ cm \ans \[TAW=0.5\times600/10\ldots=30\ \text{mm}\]
i.e.\ $0.50\times60$ cm; $RAW=0.5\times30=15$ mm \ans Supply: $RAW/ET_c=15/5=3.0$ days \finans

\exbox{Ten-soil example 1b: Sand (S03)}{Sand (S03): $\alpha=0.1450$ cm$^{-1}$, $n=2.680$, AWC $=0.6$ mm\,cm$^{-1}$. Find $m$, $S_e$ at field capacity ($h=100$ cm) and wilting point ($h=15\,849$ cm), the suction at half depletion, and TAW, RAW ($p=0.5$) and the days of supply at $ET_c=5$ mm\,d$^{-1}$ for a root zone of 0.6 m.}\sol \givenline{$\alpha=0.1450$ cm$^{-1}$ ; $n=2.680$ ; $=0.6$ mm\,cm$^{-1}$ ; $h=100$ cm ; $h=15\,849$ cm ; $p=0.5$ ; $ET_c=5$ mm\,d$^{-1}$}\ans $m=1-1/2.680=0.627$ \ans \[S_e(100)=[1+(0.1450\times100)^{2.680}]^{-0.627}=0.011\]
$S_e(15\,849)=0.0000$ \ans Half depletion is $f=0.5$, found by bisection of Eq.~(\ref{eq:fdep}): $h_{0.5}=151$ cm \ans \[TAW=0.6\times600/10\ldots=36\ \text{mm}\]
i.e.\ $0.60\times60$ cm; $RAW=0.5\times36=18$ mm \ans Supply: $RAW/ET_c=18/5=3.6$ days \finans

\exbox{Ten-soil example 1b: Loamy sand (S04)}{Loamy sand (S04): $\alpha=0.1240$ cm$^{-1}$, $n=2.280$, AWC $=0.9$ mm\,cm$^{-1}$. Find $m$, $S_e$ at field capacity ($h=100$ cm) and wilting point ($h=15\,849$ cm), the suction at half depletion, and TAW, RAW ($p=0.5$) and the days of supply at $ET_c=5$ mm\,d$^{-1}$ for a root zone of 0.6 m.}\sol \givenline{$\alpha=0.1240$ cm$^{-1}$ ; $n=2.280$ ; $=0.9$ mm\,cm$^{-1}$ ; $h=100$ cm ; $h=15\,849$ cm ; $p=0.5$ ; $ET_c=5$ mm\,d$^{-1}$}\ans $m=1-1/2.280=0.561$ \ans \[S_e(100)=[1+(0.1240\times100)^{2.280}]^{-0.561}=0.040\]
$S_e(15\,849)=0.0001$ \ans Half depletion is $f=0.5$, found by bisection of Eq.~(\ref{eq:fdep}): $h_{0.5}=172$ cm \ans \[TAW=0.9\times600/10\ldots=54\ \text{mm}\]
i.e.\ $0.90\times60$ cm; $RAW=0.5\times54=27$ mm \ans Supply: $RAW/ET_c=27/5=5.4$ days \finans

\exbox{Ten-soil example 1b: Sandy loam (S05)}{Sandy loam (S05): $\alpha=0.0750$ cm$^{-1}$, $n=1.890$, AWC $=1.2$ mm\,cm$^{-1}$. Find $m$, $S_e$ at field capacity ($h=100$ cm) and wilting point ($h=15\,849$ cm), the suction at half depletion, and TAW, RAW ($p=0.5$) and the days of supply at $ET_c=5$ mm\,d$^{-1}$ for a root zone of 0.6 m.}\sol \givenline{$\alpha=0.0750$ cm$^{-1}$ ; $n=1.890$ ; $=1.2$ mm\,cm$^{-1}$ ; $h=100$ cm ; $h=15\,849$ cm ; $p=0.5$ ; $ET_c=5$ mm\,d$^{-1}$}\ans $m=1-1/1.890=0.471$ \ans \[S_e(100)=[1+(0.0750\times100)^{1.890}]^{-0.471}=0.165\]
$S_e(15\,849)=0.0018$ \ans Half depletion is $f=0.5$, found by bisection of Eq.~(\ref{eq:fdep}): $h_{0.5}=217$ cm \ans \[TAW=1.2\times600/10\ldots=72\ \text{mm}\]
i.e.\ $1.20\times60$ cm; $RAW=0.5\times72=36$ mm \ans Supply: $RAW/ET_c=36/5=7.2$ days \finans

\exbox{Ten-soil example 1b: Loam (S06)}{Loam (S06): $\alpha=0.0360$ cm$^{-1}$, $n=1.560$, AWC $=1.7$ mm\,cm$^{-1}$. Find $m$, $S_e$ at field capacity ($h=100$ cm) and wilting point ($h=15\,849$ cm), the suction at half depletion, and TAW, RAW ($p=0.5$) and the days of supply at $ET_c=5$ mm\,d$^{-1}$ for a root zone of 0.6 m.}\sol \givenline{$\alpha=0.0360$ cm$^{-1}$ ; $n=1.560$ ; $=1.7$ mm\,cm$^{-1}$ ; $h=100$ cm ; $h=15\,849$ cm ; $p=0.5$ ; $ET_c=5$ mm\,d$^{-1}$}\ans $m=1-1/1.560=0.359$ \ans \[S_e(100)=[1+(0.0360\times100)^{1.560}]^{-0.359}=0.466\]
$S_e(15\,849)=0.0286$ \ans Half depletion is $f=0.5$, found by bisection of Eq.~(\ref{eq:fdep}): $h_{0.5}=332$ cm \ans \[TAW=1.7\times600/10\ldots=102\ \text{mm}\]
i.e.\ $1.70\times60$ cm; $RAW=0.5\times102=51$ mm \ans Supply: $RAW/ET_c=51/5=10.2$ days \finans

\exbox{Ten-soil example 1b: Silt loam (S07)}{Silt loam (S07): $\alpha=0.0200$ cm$^{-1}$, $n=1.410$, AWC $=2.0$ mm\,cm$^{-1}$. Find $m$, $S_e$ at field capacity ($h=100$ cm) and wilting point ($h=15\,849$ cm), the suction at half depletion, and TAW, RAW ($p=0.5$) and the days of supply at $ET_c=5$ mm\,d$^{-1}$ for a root zone of 0.6 m.}\sol \givenline{$\alpha=0.0200$ cm$^{-1}$ ; $n=1.410$ ; $=2.0$ mm\,cm$^{-1}$ ; $h=100$ cm ; $h=15\,849$ cm ; $p=0.5$ ; $ET_c=5$ mm\,d$^{-1}$}\ans $m=1-1/1.410=0.291$ \ans \[S_e(100)=[1+(0.0200\times100)^{1.410}]^{-0.291}=0.686\]
$S_e(15\,849)=0.0943$ \ans Half depletion is $f=0.5$, found by bisection of Eq.~(\ref{eq:fdep}): $h_{0.5}=483$ cm \ans \[TAW=2.0\times600/10\ldots=120\ \text{mm}\]
i.e.\ $2.00\times60$ cm; $RAW=0.5\times120=60$ mm \ans Supply: $RAW/ET_c=60/5=12.0$ days \finans

\exbox{Ten-soil example 1b: Sandy clay loam (S08)}{Sandy clay loam (S08): $\alpha=0.0590$ cm$^{-1}$, $n=1.480$, AWC $=1.4$ mm\,cm$^{-1}$. Find $m$, $S_e$ at field capacity ($h=100$ cm) and wilting point ($h=15\,849$ cm), the suction at half depletion, and TAW, RAW ($p=0.5$) and the days of supply at $ET_c=5$ mm\,d$^{-1}$ for a root zone of 0.6 m.}\sol \givenline{$\alpha=0.0590$ cm$^{-1}$ ; $n=1.480$ ; $=1.4$ mm\,cm$^{-1}$ ; $h=100$ cm ; $h=15\,849$ cm ; $p=0.5$ ; $ET_c=5$ mm\,d$^{-1}$}\ans $m=1-1/1.480=0.324$ \ans \[S_e(100)=[1+(0.0590\times100)^{1.480}]^{-0.324}=0.417\]
$S_e(15\,849)=0.0375$ \ans Half depletion is $f=0.5$, found by bisection of Eq.~(\ref{eq:fdep}): $h_{0.5}=369$ cm \ans \[TAW=1.4\times600/10\ldots=84\ \text{mm}\]
i.e.\ $1.40\times60$ cm; $RAW=0.5\times84=42$ mm \ans Supply: $RAW/ET_c=42/5=8.4$ days \finans

\exbox{Ten-soil example 1b: Clay loam (S09)}{Clay loam (S09): $\alpha=0.0190$ cm$^{-1}$, $n=1.310$, AWC $=1.6$ mm\,cm$^{-1}$. Find $m$, $S_e$ at field capacity ($h=100$ cm) and wilting point ($h=15\,849$ cm), the suction at half depletion, and TAW, RAW ($p=0.5$) and the days of supply at $ET_c=5$ mm\,d$^{-1}$ for a root zone of 0.6 m.}\sol \givenline{$\alpha=0.0190$ cm$^{-1}$ ; $n=1.310$ ; $=1.6$ mm\,cm$^{-1}$ ; $h=100$ cm ; $h=15\,849$ cm ; $p=0.5$ ; $ET_c=5$ mm\,d$^{-1}$}\ans $m=1-1/1.310=0.237$ \ans \[S_e(100)=[1+(0.0190\times100)^{1.310}]^{-0.237}=0.753\]
$S_e(15\,849)=0.1704$ \ans Half depletion is $f=0.5$, found by bisection of Eq.~(\ref{eq:fdep}): $h_{0.5}=618$ cm \ans \[TAW=1.6\times600/10\ldots=96\ \text{mm}\]
i.e.\ $1.60\times60$ cm; $RAW=0.5\times96=48$ mm \ans Supply: $RAW/ET_c=48/5=9.6$ days \finans

\exbox{Ten-soil example 1b: Silty clay loam (S10)}{Silty clay loam (S10): $\alpha=0.0100$ cm$^{-1}$, $n=1.230$, AWC $=1.7$ mm\,cm$^{-1}$. Find $m$, $S_e$ at field capacity ($h=100$ cm) and wilting point ($h=15\,849$ cm), the suction at half depletion, and TAW, RAW ($p=0.5$) and the days of supply at $ET_c=5$ mm\,d$^{-1}$ for a root zone of 0.6 m.}\sol \givenline{$\alpha=0.0100$ cm$^{-1}$ ; $n=1.230$ ; $=1.7$ mm\,cm$^{-1}$ ; $h=100$ cm ; $h=15\,849$ cm ; $p=0.5$ ; $ET_c=5$ mm\,d$^{-1}$}\ans $m=1-1/1.230=0.187$ \ans \[S_e(100)=[1+(0.0100\times100)^{1.230}]^{-0.187}=0.878\]
$S_e(15\,849)=0.3118$ \ans Half depletion is $f=0.5$, found by bisection of Eq.~(\ref{eq:fdep}): $h_{0.5}=906$ cm \ans \[TAW=1.7\times600/10\ldots=102\ \text{mm}\]
i.e.\ $1.70\times60$ cm; $RAW=0.5\times102=51$ mm \ans Supply: $RAW/ET_c=51/5=10.2$ days \finans

}{}
\IfFileExists{extra/tenbx01.tex}{\subsection{Worked solutions on ten crops (cool climate)}
The ten crops of the library differ in their demand, their roots and their sensitivity. The same type of question is answered for each of them at $ET_0=4.0$ mm\,d$^{-1}$; the data are typical values of the crops (Doorenbos--Kassam and FAO 56).

\exbox{Crop example 1b: Maize}{Maize grows on the loam S06 (AWC 1.7 mm\,cm$^{-1}$) with a root zone of 100 cm and a depletion fraction $p=0.55$. Find TAW, RAW and the days between irrigations at its mid-season demand $K_{c}=1.20$, $ET_0=4.0$ mm\,d$^{-1}$.}\sol \givenline{$p=0.55$ ; $K_{c}=1.20$ ; $ET_0=4.0$ mm\,d$^{-1}$}\ans \[TAW=\mathrm{AWC}\,z_r=1.7\times100=170\ \text{mm}\] \ans \[RAW=p\,TAW=0.55\times170=94\ \text{mm}\] \ans \[ET_c=K_cET_0=1.20\times4.0=4.80\ \text{mm\,d$^{-1}$}\] \ans Interval $=RAW/ET_c=94/4.80=19.5$ d \finans

\exbox{Crop example 1b: Wheat}{Wheat grows on the loam S06 (AWC 1.7 mm\,cm$^{-1}$) with a root zone of 120 cm and a depletion fraction $p=0.55$. Find TAW, RAW and the days between irrigations at its mid-season demand $K_{c}=1.15$, $ET_0=4.0$ mm\,d$^{-1}$.}\sol \givenline{$p=0.55$ ; $K_{c}=1.15$ ; $ET_0=4.0$ mm\,d$^{-1}$}\ans \[TAW=\mathrm{AWC}\,z_r=1.7\times120=204\ \text{mm}\] \ans \[RAW=p\,TAW=0.55\times204=112\ \text{mm}\] \ans \[ET_c=K_cET_0=1.15\times4.0=4.60\ \text{mm\,d$^{-1}$}\] \ans Interval $=RAW/ET_c=112/4.60=24.4$ d \finans

\exbox{Crop example 1b: Cotton}{Cotton grows on the loam S06 (AWC 1.7 mm\,cm$^{-1}$) with a root zone of 110 cm and a depletion fraction $p=0.65$. Find TAW, RAW and the days between irrigations at its mid-season demand $K_{c}=1.15$, $ET_0=4.0$ mm\,d$^{-1}$.}\sol \givenline{$p=0.65$ ; $K_{c}=1.15$ ; $ET_0=4.0$ mm\,d$^{-1}$}\ans \[TAW=\mathrm{AWC}\,z_r=1.7\times110=187\ \text{mm}\] \ans \[RAW=p\,TAW=0.65\times187=122\ \text{mm}\] \ans \[ET_c=K_cET_0=1.15\times4.0=4.60\ \text{mm\,d$^{-1}$}\] \ans Interval $=RAW/ET_c=122/4.60=26.4$ d \finans

\exbox{Crop example 1b: Tomato}{Tomato grows on the loam S06 (AWC 1.7 mm\,cm$^{-1}$) with a root zone of 70 cm and a depletion fraction $p=0.40$. Find TAW, RAW and the days between irrigations at its mid-season demand $K_{c}=1.15$, $ET_0=4.0$ mm\,d$^{-1}$.}\sol \givenline{$p=0.40$ ; $K_{c}=1.15$ ; $ET_0=4.0$ mm\,d$^{-1}$}\ans \[TAW=\mathrm{AWC}\,z_r=1.7\times70=119\ \text{mm}\] \ans \[RAW=p\,TAW=0.40\times119=48\ \text{mm}\] \ans \[ET_c=K_cET_0=1.15\times4.0=4.60\ \text{mm\,d$^{-1}$}\] \ans Interval $=RAW/ET_c=48/4.60=10.3$ d \finans

\exbox{Crop example 1b: Potato}{Potato grows on the loam S06 (AWC 1.7 mm\,cm$^{-1}$) with a root zone of 50 cm and a depletion fraction $p=0.35$. Find TAW, RAW and the days between irrigations at its mid-season demand $K_{c}=1.15$, $ET_0=4.0$ mm\,d$^{-1}$.}\sol \givenline{$p=0.35$ ; $K_{c}=1.15$ ; $ET_0=4.0$ mm\,d$^{-1}$}\ans \[TAW=\mathrm{AWC}\,z_r=1.7\times50=85\ \text{mm}\] \ans \[RAW=p\,TAW=0.35\times85=30\ \text{mm}\] \ans \[ET_c=K_cET_0=1.15\times4.0=4.60\ \text{mm\,d$^{-1}$}\] \ans Interval $=RAW/ET_c=30/4.60=6.5$ d \finans

\exbox{Crop example 1b: Soybean}{Soybean grows on the loam S06 (AWC 1.7 mm\,cm$^{-1}$) with a root zone of 60 cm and a depletion fraction $p=0.50$. Find TAW, RAW and the days between irrigations at its mid-season demand $K_{c}=1.15$, $ET_0=4.0$ mm\,d$^{-1}$.}\sol \givenline{$p=0.50$ ; $K_{c}=1.15$ ; $ET_0=4.0$ mm\,d$^{-1}$}\ans \[TAW=\mathrm{AWC}\,z_r=1.7\times60=102\ \text{mm}\] \ans \[RAW=p\,TAW=0.50\times102=51\ \text{mm}\] \ans \[ET_c=K_cET_0=1.15\times4.0=4.60\ \text{mm\,d$^{-1}$}\] \ans Interval $=RAW/ET_c=51/4.60=11.1$ d \finans

\exbox{Crop example 1b: Sugar beet}{Sugar beet grows on the loam S06 (AWC 1.7 mm\,cm$^{-1}$) with a root zone of 90 cm and a depletion fraction $p=0.55$. Find TAW, RAW and the days between irrigations at its mid-season demand $K_{c}=1.20$, $ET_0=4.0$ mm\,d$^{-1}$.}\sol \givenline{$p=0.55$ ; $K_{c}=1.20$ ; $ET_0=4.0$ mm\,d$^{-1}$}\ans \[TAW=\mathrm{AWC}\,z_r=1.7\times90=153\ \text{mm}\] \ans \[RAW=p\,TAW=0.55\times153=84\ \text{mm}\] \ans \[ET_c=K_cET_0=1.20\times4.0=4.80\ \text{mm\,d$^{-1}$}\] \ans Interval $=RAW/ET_c=84/4.80=17.5$ d \finans

\exbox{Crop example 1b: Sunflower}{Sunflower grows on the loam S06 (AWC 1.7 mm\,cm$^{-1}$) with a root zone of 80 cm and a depletion fraction $p=0.45$. Find TAW, RAW and the days between irrigations at its mid-season demand $K_{c}=1.10$, $ET_0=4.0$ mm\,d$^{-1}$.}\sol \givenline{$p=0.45$ ; $K_{c}=1.10$ ; $ET_0=4.0$ mm\,d$^{-1}$}\ans \[TAW=\mathrm{AWC}\,z_r=1.7\times80=136\ \text{mm}\] \ans \[RAW=p\,TAW=0.45\times136=61\ \text{mm}\] \ans \[ET_c=K_cET_0=1.10\times4.0=4.40\ \text{mm\,d$^{-1}$}\] \ans Interval $=RAW/ET_c=61/4.40=13.9$ d \finans

\exbox{Crop example 1b: Onion}{Onion grows on the loam S06 (AWC 1.7 mm\,cm$^{-1}$) with a root zone of 35 cm and a depletion fraction $p=0.30$. Find TAW, RAW and the days between irrigations at its mid-season demand $K_{c}=1.05$, $ET_0=4.0$ mm\,d$^{-1}$.}\sol \givenline{$p=0.30$ ; $K_{c}=1.05$ ; $ET_0=4.0$ mm\,d$^{-1}$}\ans \[TAW=\mathrm{AWC}\,z_r=1.7\times35=60\ \text{mm}\] \ans \[RAW=p\,TAW=0.30\times60=18\ \text{mm}\] \ans \[ET_c=K_cET_0=1.05\times4.0=4.20\ \text{mm\,d$^{-1}$}\] \ans Interval $=RAW/ET_c=18/4.20=4.2$ d \finans

\exbox{Crop example 1b: Alfalfa}{Alfalfa grows on the loam S06 (AWC 1.7 mm\,cm$^{-1}$) with a root zone of 100 cm and a depletion fraction $p=0.55$. Find TAW, RAW and the days between irrigations at its mid-season demand $K_{c}=1.15$, $ET_0=4.0$ mm\,d$^{-1}$.}\sol \givenline{$p=0.55$ ; $K_{c}=1.15$ ; $ET_0=4.0$ mm\,d$^{-1}$}\ans \[TAW=\mathrm{AWC}\,z_r=1.7\times100=170\ \text{mm}\] \ans \[RAW=p\,TAW=0.55\times170=94\ \text{mm}\] \ans \[ET_c=K_cET_0=1.15\times4.0=4.60\ \text{mm\,d$^{-1}$}\] \ans Interval $=RAW/ET_c=94/4.60=20.3$ d \finans

}{}
\IfFileExists{extra/ex01.tex}{\backhead{EXERCISES}{khred}
\begin{exlist}
\item A soil core of volume 200 cm$^3$ has a dry mass of 270 g. Find the bulk density, the porosity ($\rho_s=2.65$) and the void ratio.\exans{1.35 g cm$^{-3}$; 0.491; 0.963}
\item The gravimetric water content of a soil with $\rho_b=1.42$ g\,cm$^{-3}$ is 0.18. Find $\theta$, the degree of saturation and the depth of water in a 60 cm layer.\exans{0.256; 0.55; 153 mm}
\item A 1.2 m profile contains 312 mm of water when its mean content is measured. By how many millimetres must it be raised to reach $\theta=0.34$? What volume of water in m$^3$ does this need for 5 ha?\exans{96 mm; 4800 m$^3$}
\item Tensiometers at 20 cm and 50 cm depth read 180 cm and 140 cm. Find the heads (datum at the surface), the direction of flow and the gradient.\exans{$-200$ and $-190$ cm; upward; 0.333 cm per cm}
\item The loam S06 has $\alpha=0.036$, $n=1.56$. Find $S_e$ at suctions of 50, 300 and 1000 cm.\exans{0.638, 0.262, 0.134}
\item For the sandy loam S05 ($\alpha=0.075$, $n=1.89$, $\theta_s=0.41$, $\theta_r=0.03$) find the plant-available water per metre between 100 and 15\,849 cm.\exans{62 mm per m}
\item A root zone 0.6 m deep has TAW = 96 mm and $p=0.4$. What is the readily available water? How long does it last at $\mathrm{ET}_c=4.8$ mm\,d$^{-1}$?\exans{38.4 mm; 8.0 d}
\item A drip system applies 18 mm per event to a field of 3.5 ha. What volume of water is needed? How long does the event last if the pump delivers 25 m$^3$\,h$^{-1}$?\exans{630 m$^3$; 25.2 h}
\item The water balance of a rainfed year gives $P=520$, $T=230$, $E_s=180$, $E_o=40$, $D=60$, $Q=10$ mm. Find $\Delta S$. If $I=120$ mm were added and all went to $T$, what would the new $T$ be?\exans{$0$ mm; 350 mm}
\item Capillary rise in a tube of radius $r=0.01$ mm with $\sigma=0.0728$ N\,m$^{-1}$, $\gamma=0$: find $h$ (cm).\exans{148 cm}
\end{exlist}
}{}
\IfFileExists{extra/dyk01.tex}{\dykband{DO YOU KNOW ?}
\begin{exlist}
\item Define porosity, void ratio and bulk density. How are they related?
\item Distinguish gravimetric from volumetric water content. Why does a model use the latter?
\item Name the four components of the total soil-water potential. Which ones drive flow in the soil, and which acts only on root uptake?
\item What are field capacity and the permanent wilting point in terms of suction? Why are these values conventions and not constants of the soil?
\item Why does the van Genuchten parameter $n$ matter more than $\alpha$ for the plant-available water?
\item Why is a millimetre of water over a hectare equal to 10 m$^3$?
\item Describe three methods of measuring soil water and state what each one integrates.
\item Why is the conductivity of a soil not a constant?
\end{exlist}
}{}
\IfFileExists{extra/obj01.tex}{\dykband{OBJECTIVE TYPE QUESTIONS}
\begin{exlist}
\item The dry bulk density of a soil is $1.325$ g\,cm$^{-3}$; its porosity (with $\rho_s=2.65$) is\opt{0.25}{0.50}{0.75}{1.00}
\item The volumetric water content equals the gravimetric content multiplied by\opt{$\rho_s/\rho_w$}{$\rho_b/\rho_w$}{$\rho_w/\rho_b$}{$\phi$}
\item The suction at field capacity used in this book is\opt{15 cm}{100 cm}{330 cm}{15\,849 cm}
\item The pF of the wilting point is\opt{2.0}{3.0}{4.2}{7.0}
\item A large value of the van Genuchten $n$ indicates\opt{a wide pore-size distribution}{a uniform pore size}{a clay soil}{a high organic matter}
\item Water in an unsaturated soil moves\opt{from low to high head}{from high to low head}{always downward}{towards high suction only}
\item 1 mm of water over 1 ha is\opt{1 m$^3$}{10 m$^3$}{100 m$^3$}{1000 m$^3$}
\item The osmotic potential affects\opt{flow in the soil only}{uptake by roots}{the bulk density}{the porosity}
\item Which method gives the water content by weighing?\opt{TDR}{neutron probe}{oven drying}{tensiometer}
\item Readily available water is TAW multiplied by\opt{the root depth}{the depletion fraction $p$}{the efficiency}{the porosity}
\end{exlist}
\objkey{1.~(b)\quad 2.~(b)\quad 3.~(b)\quad 4.~(c)\quad 5.~(b)\quad 6.~(b)\quad 7.~(b)\quad 8.~(b)\quad 9.~(c)\quad 10.~(b)}
}{}
